Les entretiens qui suivent ont été réalisés par Mélina Franot au cours d’un stage de 6 mois à la Direction des Patrimoine. Toutes les personnes ayant participé à ces entretiens ont accepté d’apparaitre dans ce mémoire, dont une anonymement. Les entretiens enregistrés puis transcrits ont subi quelques légères modifications (suppression de répétitions, d’hésitation, tics de langages…) afin d’améliorer le confort de lecture et la compréhension. Conformément au Règlement Général sur la Protection des Données (RGPD) les acteurs interrogés ont eu droit d’accès et de rectification.\newline

Ces échanges prenant la forme d’entretiens dirigés, certaines questions ont été posées à chaque personnes et certaines à une seule, selon le déroulement de l’entretien.

\section*{Présentation des participants à ces entretiens}

\subsection*{Documentaliste}  

\textbf{Documentaliste :} […] l'INA a organisé un concours de recrutement en vue de la mise en œuvre du dépôt légal de la radio-télévision.  On était en 1993. Mais le dépôt légal n'a pas été mis en œuvre tout de suite, malheureusement, parce qu'il y a eu un changement de majorité politique et le décret d'application qui devait être passé a été retardé d'un an et demi. J'ai eu la chance entre-temps que l'INA me propose un CDD à la Vidéothèque d'Actualité en attendant la mise en œuvre effective. Parce que, qui dit décret d'application dit aussi budget. L’INA avait organisé un concours de recrutement mais ne pouvait plus proposer les postes prévus. J’étais donc au service des “commandes” à la Vidéothèque d'actualité, rue de Patay à Paris. Je répondais en priorité aux commandes de TF1, mais aussi aux commandes dites “commerciales” (émanant d’autres producteurs, d'autres chaînes privées françaises ou des chaînes étrangères). À l'époque, les occasionnels n'avaient pas le droit de toucher à l'indexation puisqu'ils n'étaient pas formés. Il n'existait pas encore de formation de documentaliste audiovisuel. On était formés après le concours. Les débutants étaient jetés “dans le grand bain” des commandes, ce qui n'était pas simple non plus. J'ai fait ça pendant un an. Ensuite, j'ai travaillé à la préparation de la saisie des fichiers de la Vidéothèque d’actualité. Avant 1975 et l'introduction d'un système informatique à l’INA, les émissions étaient décrites sur des notices cartonnées classées dans des tiroirs. Ces tiroirs étaient en cours de saisie dans la base informatique et j'étais chargée, en tant que CDD, avec des collègues CDI planifiés à tour de rôle, de préparer ces tiroirs pour compléter les données avant de les envoyer à un prestataire. Au bout d'un moment, les tiroirs ont été envoyés directement au prestataire pour compléter les données ensuite. Mais là, on était dans une première phase de test. Il s'agissait de compléter les noms des journalistes, de vérifier la présence du matériel, etc., des choses comme ça.  Cela, jusqu'à la création du dépôt légal. Et donc fin 1994, toutes les personnes qui avaient été recrutées lors du concours un an et demi avant, plus d'autres qui avaient été repêchées - parce que certaines personnes avaient dû se désister, avec déjà un emploi ou ne pouvant pas attendre. Donc, on a commencé la mise en œuvre du dépôt légal, officiellement le 1er janvier 1995.  On a eu trois mois de formation.  Pas une formation uniquement sur l'indexation, aussi sur l’histoire de la télévision, de la radio, la sémiologie, le droit de l’information, avec une visite des services de l'INA. Et des modules de travaux pratiques, où l’on nous expliquait comment analyser les émissions, faire les résumés, travailler avec le thésaurus, utiliser les logiciels développés pour ce nouveau département, etc. A l'issue des trois mois de formation, en février 1995, on a commencé à travailler pour le dépôt légal. Je suis restée au service du dépôt légal jusqu'à la réorganisation récente des départements de l'INA, qui a éclaté un peu tous les services. Je fais partie de ce qu'on appelle les ex-DL, qui y ont fait leur parcours pendant plusieurs années, sauf au début où j'ai travaillé à la Vidéothèque d'actualité, que l’on appelait “les archives professionnelles” (à destination des “professionnels”, journalistes, producteurs). Ensuite, je suis restée au dépôt légal, que j'aimais bien, parce qu'on travaillait à la fois sur la radio, la télévision, et sur plusieurs chaînes. J'aimais bien indexer, faire des résumés.

\subsection*{Pascal Lebeurier}

\textbf{PL :} Je suis à L’INA depuis 30 ans. J'ai été documentaliste de 1995 à 2012, plutôt côté actualité. À l'époque, il y avait séparation entre actualité, production, phonothèque et j'étais plutôt actualité. Je suis devenu cadre en documentaire en 2012, c'est la partie commande. Donc, j'encadre les documentalistes qui répondent aux demandes clients, ça peut être de la demande pour un JT : TF1, comme pour un documentaire : l'histoire de l'Algérie qu'on avait fait il y a 3-4 ans, en 5 fois 52 minutes, ça peut être des formats très différents. Donc les recherches sont un peu différentes quand tu travailles pour une expo à cinémathèque, ou pour un documentaire Arte, tu travailles pour un JT TF1, c'est assez différent. J'ai aussi refait un master histoire et patrimoine en 2018 et ça m'a donné du recul sur l'archivistique et la prise de recul sur les demandes, notamment histoire. 

\textbf{MF :} Quelles sont vos missions actuelles ?  

\textbf{PL :} Je suis toujours responsable de la partie commande, on appelle ça plutôt communication des fonds.  Des fonds, mais plutôt aux clients professionnels, pas aux usages du publics, il y a des services à part, moi, c'est Clients et Partenaires. Ce sont les JT privés comme TF1, CNews ; tous les documentaires faits par l'INA, les productions : ce sont des émissions INA comme ADN, Rambobina. Là, il y a une nouvelle émission qui s'appelle « Témoins d'Actu », donc des émissions INA. C'est aussi la partie commerciale, les projets commerciaux. Donc ça peut être des plateformes, comme Netflix, ça peut être plein de choses. Et puis il y a la partie SAC, donc c'est la partie culturelle, c'est-à-dire les expos, à la Cinémathèque, à la mairie de Paris.  

\subsection*{Arthur Hénaff}

\textbf{AH :} Je suis Arthur Hénaff, je suis responsable de l'INAthèque depuis octobre 2025 et je fais partie de l'équipe de pilotage du Lab de l'INA. Avant ça, je suis conservateur de bibliothèque en détachement à l'INA. L'INAthèque, pour le dire très rapidement, c'est un des services du département de valorisation scientifique à la direction des patrimoines. L'autre service étant « manifestation et partenariat scientifique ». Par ailleurs, il y a aussi les éditions de l'INA qui se trouvent dans ce département.

\section*{Bases de données de l’INA }

\subsection*{Documentaliste} 

\textbf{MF :} Et du coup, là, actuellement, au niveau des logiciels, des bases de données surtout, lesquelles tu utilises ? 

\textbf{Doc :} En fait, maintenant, je travaille sur tous les logiciels, parce qu'en étant à TEC, Traitement et Enrichissement des Collections, je continue le traitement documentaire du flux, donc j'utilise les anciens outils du dépôt légal, qui s'appellent DLTV, DL Radio. Ce sont les outils de saisie qui travaillent avec Oracle. On utilise le logiciel de consultation des données Hyperbase. En même temps, je fais aussi de la “reprise d'antériorité”, c'est-à-dire du traitement documentaire d'anciennes collections du secteur professionnel, enfin ex-professionnel, et donc sur Totem. Je travaille aussi sur les inventaires, ce qu'on appelle les dépôts de fonds, ce sont des fonds audiovisuels, que des personnalités de la radio ou de la télévision déposent à l'INA, ou leurs descendants, ou des institutions. C'est Malika Mekdad qui s'occupe de ça, et je participe à certains inventaires. Là, on crée dans Totem des notices qui seront reliées au matériel déposé. Totem, je ne le connais pas de façon aussi fine que les collègues qui travaillent pour les archives professionnelles, qui y font des recherches. Je ne suis pas aussi experte, mais je me débrouille avec ce dont j'ai besoin.  Mais je sais qu'il y a des fonctionnalités que je ne connais pas très bien dans Totem.

\textbf{MF :} Mais du coup, même en l'utilisant pas forcément à 100\%, est-ce que tu vois des incohérences dans les données, qui t'empêchent en tout cas d'avancer dans ce que tu fais ?  

\textbf{Doc :} Dans les données ? Par rapport à ce qu'on voit ou par rapport à la façon de rechercher ?  

\textbf{MF :} C'était plus vis-à-vis de ce que l'on voit, mais c'est vrai que dans la façon de chercher, c'est intéressant aussi. 

\textbf{Doc :} Les stratégies de recherche sont différentes entre Hyperbase et Totem. De fait, mes recherches, par habitude, je les fais plutôt dans Hyperbase parce que je me sens plus à l'aise dans cette logique de recherche qui combine des requêtes. Les usages étaient différents.  Hyperbase a été conçu pour créer des corpus, pour des chercheurs, des étudiants. Donc en fait, pour agglomérer des recherches, des requêtes, et constituer peu à peu un corpus. Par association d'idées, recombiner les requêtes pour élargir la recherche, ou faire des petites variations pour récupérer des notices qui nous auraient échappé, dans le but de créer un corpus, en tournant la base dans tous les sens. Totem, c'est plus à vocation des professionnels, donc pour des recherches d'images ou de sons précis, généralement. Il a été constitué plutôt pour faire des recherches puis les affiner. Totem est un peu postérieur à Hyperbase avec des fonctionnalités plus récentes. Ils avaient pu créer ce qu'on appelle des facettes. À partir d'un lot résultat, l’outil analyse ce lot et permet d'éliminer certains éléments dont on ne veut pas, sélectionner une chaîne, par exemple, ou une collection. Ça permet d'affiner la recherche. On peut le faire dans HyperBase en rajoutant une requête. C'est une autre approche, une autre logique. Dans Totem, on affine la recherche pour pouvoir cibler des émissions pertinentes que le client veut retrouver. C'est une approche un peu différente, une logique un peu différente. Dans Totem, on peut plus difficilement combiner des requêtes. On peut le faire, mais c'est un peu plus laborieux, disons qu'il faut d'abord avoir un lot résultat et puis après le récupérer pour le combiner à une autre requête. Alors que dans Hyperbase, on peut le faire très facilement en deux secondes. Ce sont des usages différents dès l'origine, donc des conceptions aussi un peu différentes. Et puis selon les habitudes que l'on a... J'ai des collègues côté archives professionnelles qui préfèrent utiliser Totem parce que, pour eux, ils sont plus familiers de ce module-là. Et inversement, Hyperbase, ils trouvent ça trop compliqué. Donc chacun fait selon ses habitudes.

\textbf{MF :} Mais du coup, même s'il n'y a plus clairement et institutionnellement la différence entre pro et DL, les habitudes font que 

\textbf{Doc :} Dans Totem, il y a des émissions, des notices qui sont présentes et qui ne seront pas dans Hyperbase, pour différentes raisons.  Dans ces cas-là, je commence toujours mes recherches dans Hyperbase, mais après je complète dans Totem pour être sûre de ne pas manquer quelque chose, parce que je sais que tout ne sera pas dans Hyperbase, où l’idée est de rendre compte de ce qui était diffusé, donc c'est un peu la priorité. Dans Totem, il y aura d’autres choses, comme des éléments de montage, des éléments de travail, des choses comme ça, qui ne seront pas dans Hyperbase, parce que ce n'est pas diffusé, ou sont à usage interne à l’INA.  La mission du dépôt légal, c'est de rendre compte de ce qui est diffusé, c'est la première mission. Alors que dans Totem, il y a toutes sortes de documents dont les étudiants, enfin le public, l'extérieur, n'a rien à faire. En fait, les deux outils se complètent. Après, il y a le fameux Notilus, un nouvel outil, qui réunit tous les fonds. J’ai un peu de mal à faire des recherches dans Notilus donc je préfère travailler avec ce que je sais. Je sais où je mets les pieds, je vois un petit peu si ça répond ou si ça répond pas je peux savoir si j'ai fait une erreur dans ma question, si je peux rebondir de telle ou telle façon. Dans Notilus, je me sens moins à l'aise. Mais par rapport à ce que tu disais, ce qu'on trouve dans les notices, justement, c'est le fait de ces périmètres un peu différents entre Totem et Hyperbase.  C'est pour ça que là, c'est intéressant. Moi, je sais qu'il ne faut pas se contenter de l'un ou l'autre selon les types de recherches que je dois faire.  Parce que je sais qu'il y a des éléments qui figureront dans Totem et que je n'aurai pas dans Hyperbase. C'est difficile à décrire comme ça, mais il y a des choses qu'on n'aura pas de la même façon. C'est difficile à expliquer. Il y a par exemple des champs qui n'apparaissent pas du tout dans Hyperbase. Il y a le champ « dispositif » dans la notice qui donne par exemple le plan par plan de certaines émissions.  Eh bien ça ne s'affichera pas dans Hyperbase. 

\textbf{MF :} D'accord. Pour une même émission ? 

\textbf{Doc :} Absolument. Ça pose question, par exemple pour France 2. On ne traite plus le journal de 20h depuis quelques temps. Or, France 2 fait une description plan par plan, qui arrive dans le champ séquence. Et les chercheurs n'ont pas accès à ce champ dans Hyperbase. Donc, pour faire des recherches... Il y a des petites choses comme ça que l'on sait et dont on doit tenir compte quand on fait des recherches particulières.

\textbf{MF :} Comme ce sont des logiciels qui sont quand même vieillissants, est-ce ça peut poser problème ? Est-ce qu'ils peuvent planter ou ne pas fonctionner comme tu le voudrais ? 

\textbf{Doc :} Ça, non.  On n'est pas confronté à ça.  Les seuls plantages qui pourraient arriver, ce sont les plantages de serveurs.  Oui, dans ces cas-là, il n'y a plus rien.  Enfin, je veux dire, on n'accède plus à rien.  Mais des bugs d'affichage ou des choses bizarres, non.  Moi, j'ai une confiance.  S'il y a eu des bugs, c'était dans les débuts, mais ça s'est stabilisé avec le temps. Donc, je sais que si je n'ai pas de résultat, c'est plutôt moi. Je m'interroge sur ma stratégie de recherche.  Je me dis que je n'ai peut-être pas utilisé le bon mot, le bon vocabulaire. Mais c'est lié à l'expérience aussi de l'outil.  

\subsection*{Pascal Lebeurier}

\textbf{MF :} Du coup de votre côté, c'est la base de données totem ? 

\textbf{PL :} Nous c'est totem et un peu Notilus maintenant. Comme Notilus fait des choses que Totem ne fait pas notamment la recherche dans les transcriptions.  

\textbf{ML :} Quand vous n'utilisiez que Totem, puisque Nautilus, c'est encore assez récent, est-ce qu'il pouvait vous manquer certaines informations sans la transcription ? 



\textbf{PL :} Oui, en fait, c'est plutôt le problème de l'histoire de l'indexation des fonds, on a récupéré la cinémathèque à l’ORTF, on a un fonds tellement énorme que parfois, il y a juste un titre, parfois, il manque des infos. Donc ça veut dire que si tu cherches un nom, il peut ne pas être dans la notice. Donc la transcription, ça va nous aider par rapport à ça.  Si tu cherches, je ne sais pas, Charles Aznavour, tu peux très bien essayer de chercher dans la transcription, s'il n'y a pas d'émissions qui ne sont pas indexées, dans lesquelles il y a Charles Aznavour qui y est. Mais autrement, on va chercher. Les documentalistes, en plus, sont expérimentés, certains connaissent très bien les fonds, et donc ça veut dire qu'ils savent qu'il faut aller chercher à côté, si tu veux trouver. Tu vas d'abord faire une recherche avec le nom, au générique, tu vas avoir déjà beaucoup de résultats, mais tu sais que tu peux aller chercher ailleurs pour trouver d'autres émissions. 

\section*{Impact de la politique documentaire sur les données de description }

\subsection*{Documentaliste} 

\textbf{MF :} Donc ça arrive parfois que vous n’ayez pas forcément la réponse pour les chercheurs ? 

\textbf{Doc :} Ça, ça arrive, oui. Je sais qu'à un moment donné, par exemple, c'étaient les journaux de « Soir 3 », je pense que ça a été rattrapé depuis, mais à un moment donné, on ne les indexait plus pour x raison, France 3 ne payait plus l’Ina pour ça, ou quelque chose comme ça. Bon, j'étais jeune CDD à l'époque, donc prudence avec ce que je raconte. Mais ça, c'était à l'époque où il n'y avait pas encore vraiment le dépôt légal qui était mis en route. Donc si on voulait faire des recherches sur ces fonds-là, il n'y avait rien qui ressortait parce que ce n'était pas décrit. Et les jours de grève... Parfois, c'est marqué « non indexé parce qu'en grève ». Donc, ça arrive. D'où le rôle de la reprise d’antériorité. Ca fait partie de ces campagnes qui permettent de revenir sur des collections qui ont été peu ou partiellement décrites, enfin, partiellement décrites pour X raisons. Soit manque de temps, manque de moyens, soit que grève, soit que machin. Parce qu'il y a toujours eu des problèmes de temps, de recrutement, des choses comme ça.  Un exemple, c’est « Aujourd'hui Madame ». On a entrepris un gros travail de reprise d’antériorité pendant deux ans, avant l'arrivée de la transcription automatique. Ceci dit, ça n'aurait peut-être pas résolu la question. C'est une collection qui a été diffusée entre 1969-1982, 2000 émissions environ, cinq émissions par semaine. Et elle n'avait pas forcément été traitée de façon très fine tout le temps, pour différentes raisons. On s'est rendu compte que cette émission avait un intérêt à la fois historique et sociologique, parce qu’elle abordait diverses questions de société. C’étaient beaucoup des débats en plateau, avec des téléspectatrices invitées à participer. Il y avait aussi des hommes, mais, comme c'était une émission de l'après-midi, ça a touché plus les femmes. Elle a abordé toutes sortes de sujets de société et, évidemment, le temps passant, ça éveille des échos avec notre temps. Les personnes qui étaient aux “commandes” pour les professionnels, se rendaient compte de l'intérêt de cette émission et de l'intérêt de la décrire davantage, pour ne pas passer à côté. C’étaient essentiellement des discussions en plateau, mais il y avait aussi des micros-trottoirs auprès de jeunes, d'enfants, de vieux, sur tous les sujets, des reportages, etc. C'était vraiment précieux. Pour la reprise, on a été planifié à tour de rôle. Ce n'est pas la même personne qui a fait 2000 émissions. Pour reprendre toutes les notices, compléter celles qui n'étaient pas décrites ou d'autres partiellement, parce que parfois il y avait juste une partie de l'émission qui était décrite. L’autre partie n'était peut-être pas jugée à l’époque intéressante pour une reprise, pour une commercialisation ultérieure. Il y avait ça aussi, il y avait ce regard professionnel d'intérêt commercial ou pas. On n'avait pas de notion de patrimoine audiovisuel à l'époque, pas encore. C'est venu plus tard. Donc, c'est ça qui faisait aussi que dans les descriptions, il y avait parfois uniquement une toute petite partie de l'émission qui était décrite.  Et le reste, “pas intéressant, on s'en fiche”. Alors qu'aujourd'hui, on dit « Ah, mais si ! ».

\textbf{MF :} Oui, donc il y a tout une évolution, justement, de la politique documentaire ? 

\textbf{Doc :} Absolument.  Le regard sur l'archive, tout simplement. C'était l’essor de la télévision, mais les universitaires ne considéraient pas encore que la télévision ou la radio étaient un objet légitime d’étude, comme une source de recherche. Il y avait quelques pionniers, mais c'était vraiment à la marge. Ce regard-là, c'est venu plutôt dans les années 80. On était dans un usage uniquement commercial, professionnel. Et puis le regard sur les sujets de société a évolué.  Évidemment, il y a des questions qui, à l'époque, peut-être, n'apparaissaient pas forcément intéressantes, alors qu'aujourd'hui, on se dit, « ben si ! », au contraire. L'autre question par rapport à la pertinence des notices, il y a une grave question, mais qui n'a malheureusement pas pu être vraiment résolue, c'est qu'en fait, on n'a pas gardé les notices d'origine. En tout cas, pour la partie après 75, parce qu'avant 75, si c'est sur des notices papiers, celles-là, elles ont été conservées. Mais après 75, on a réécrit par-dessus. Donc, on n'a pas la trace du vocabulaire. Et ça, c'est vrai que c'est un peu gênant, le côté patrimonial. Mais en même temps, le logiciel ne permet pas non plus de...  C'est-à-dire que dans le champ notes, on met quand même que la collection a fait l'objet d'un traitement documentaire à telle date.  Pour dire que voilà, ça a été divisé en 73, mai 75, mais ça a été revu et corrigé en 2021-22. On sait que ça a été retrafiqué. Et puis sans parler que toutes les notices, même avant, par exemple, les collègues, quand ils tombent sur une notice qui était avant en papier, il y avait juste le titre et puis un plan par plan, quand ils tombent dessus, souvent ils rajoutent les mots-clés et puis rédigent un petit peu parfois quand ils ont le temps donc ça aussi ça a été refait et corrigé donc la notice d'avant a disparu. 

\textbf{MF :} Du coup dans ce cas-là c'est revu systématiquement ou c'est vraiment quand on passe au cas par cas ? 

\textbf{RD :} En fait c'est au cas par cas. Il y a eu des campagnes de reprise d'antériorité mais, pour les sujets de journaux télévisés, c'est plutôt au cas par cas. En fait, ça dépend un peu des thématiques. Par exemple, pour les périodes comme la guerre d'Algérie, ça a été fait de façon systématique, parce que c'est tellement demandé, tellement précieux, que là, pour le coup, c'est bien balisé. On sait qu'il y a des sujets, pour les fonds anciens, on n'a pas besoin de se poser trop de questions, on sait que c'est traité, c'est carré. Le fonds de presse filmée des Actualités françaises aussi, ça a été bien balisé, tout ça.  Mais on sait qu'il y a beaucoup de notices qui ne sont vraiment pas… Parce que ce sont des sujets un peu pointus ou autre et du coup ça n'a pas été complété.  Il y a tous les fonds régionaux aussi qui sont partiellement, très succinctement décrits parce qu'on s'attachait plutôt aux nationales jusqu'à une date récente. Maintenant ça devient très précieux. D'où l'intérêt des nouveaux outils, parce que la transcription automatique va permettre de récupérer sans doute, de repêcher des masses de documents sur lesquels on ne pouvait pas passer de toutes façons. Ça va être fantastique de pouvoir récupérer tout ça parce que, parfois, on peut avoir des entretiens avec des artistes, lorsqu’il y a des expositions, enfin, tout ce genre de choses.  Donc, on va peut-être faire des découvertes.

\textbf{MF :} C'est quelque chose qui est prévu ou qui serait idéalement ?   

\textbf{RD :} Je pense que ça va se faire, avec la transcription automatique. À mon avis, c'est une question de mois ou d'années. Enfin, moi, je ne suis pas dans les petits papiers, mais j'imagine que ça va forcément venir. Du coup, le travail de reprise d'antériorité tel qu'on l'a conçu jusqu'ici, ça va être amené sans doute à être adapté. On va peut-être prioriser sur certains types de...  Enfin, c'était déjà priorisé, parce que de toute façon, on a ciblé certaines collections jugées intéressantes. J'ai fait par exemple, à un moment donné, un tableau sur les émissions jeunesse, où j'étais amenée à prioriser : telle collection, ça serait bien qu'on puisse la décrire.  Il y a par exemple une émission, ce n'est pas moi qui l'ai détectée, parce qu'elle était bien connue auparavant, qui s'appelait « Pile ou Face », dans les années 50, une sorte de leçon de choses. Il y avait des reportages sur toutes sortes de sujets, les métiers, etc. A l'époque, comme c'était une émission pour la jeunesse, elle n'était quasiment pas décrite. Mais, avec le recul des années, on s'est rendu compte que ça devenait une source précieuse, parce qu'on a des plans sur des artisans, parfois des métiers qui ont disparu, des choses comme ça. En plus, souvent très bien filmé, un beau noir et blanc. C'est le genre de collection qu'il est précieux de revoir et de compléter. Et là, pour le coup, la transcription automatique n'est pas d'une grande aide, parce que c’est souvent sans bande son. 

\subsection*{Pascal Lebeurier  }

\textbf{MF :} Oui, ce n'est pas la même politique des archives (aux Etats-Unis). 

\textbf{PL :} Oui, ils n'ont pas la même les Américains. Après, moi, je ne suis pas spécialiste.  Il faudrait interroger Serge Viallet ou Dominique Forger, qui sont vraiment des spécialistes.  Ils allaient tous les ans là-bas. Mais en fait, ils expliquent que là-bas, tu fais des recherches, ils te sortent les supports, tu as aussi tous les documents d'accompagnement. Serge Viallet m’avait montré tout ce qu'il avait ramené sur le débarquement. Il a une tonne de documentation écrite sur comment les Américains ont filmé le débarquement.  

\textbf{MF :} Oui, c'est un traitement des archives très différents. 

\textbf{PL :} Oui, c'est ça. 

\textbf{MF :}  Et puis l’INA se place quand même un peu, comme une figure de proue des archives audiovisuelles. 

\textbf{PL :} Oui, parce qu'en fait, c'est assez unique l'INA, le fait qu'on ait autant de fonds, qu'on soit un organisme qui gère différentes chaînes. En Italie, tu vas avoir « Luce », qui est le fonds archive-film, et après, tu vas avoir la RAI. Peut-être que la BBC a un gros fonds aussi. Mais c'est vrai que dans d'autres pays, il n'y a pas l'INA. Donc on a un volume de fonds qui est peut-être unique, en effet. Mais qui dit volume de fond, c'est ce que je dis souvent, je fais des cours aux étudiants, au master patrimoine. Le problème, c'est que, le volume de fond est important, donc, si tu ne sais pas chercher, tu vas avoir énormément de bruit. Donc, c'est pour ça que nous, des fois, on nous demande des choses, pour ina.fr, on va te demander en une demi-heure. Bon, pour France Télé, quand les docs sont à France Télé, ils vont te dire, en 30 minutes, fais-moi une sélection de 10 documents sur Édith Cresson ou Edouard Balladur. Pour faire ça rapidement, il faut être expérimenté, savoir trouver 10 documents qui retracent la carrière d'Édouard Balladur. Tu peux t'aider de la doc écrite, tu peux te dire, il a été ministre, il a été premier ministre, et après tu te sors un document. Tu vas mettre sa nomination au premier ministre, tu vas faire des choses comme ça. Tu t'aides aussi de la doc écrite pour faire la recherche dans les fonds. Mais nous, on a des fonds tellement importants, qu'en effet, il faut avoir des méthodes, autrement, tu vas être noyé. Autrement, au bout d'une demi-heure, tu vas dire, ah ben non, laissez-moi. Et le journaliste, il va dire, « ben non, le 13H, il est à 13h, moi, il me faut les trucs tout de suite ».  Alors après, il faut s'habituer, tu dis que ce que tu proposes, c'est une partie, tu vas pas être exhaustif, tu vas peut-être faire un raté, mais pour autant qu'il y ait du contenu pour faire une rétro, voilà.  Plus t'as de fond, plus faut être dans la recherche, être spécialisé, avoir une expérience de la recherche, ça c'est important. C'est pour ça que sur ina mediapro, les personnes qui cherchent eux-mêmes, parfois elles se noient. Elles se noient et au bout d'un moment elles reviennent vers l'INA en disant « Désolé, mais je n'y arrive pas ».  Ça fait cinq jours que je travaille et je n'en suis qu'à la première année. Mais par exemple, sur les AF, quand je parle d'archivistique, c'est vrai que moi, j'avais poussé pour qu'on mette une note. Sur ina mediapro, il y a des personnes qui tombaient sur les activités allemandes ou Vichy, sans se dire que c'était des actualités propagandes. Maintenant on a mis une note. Comme en archivistique, tu auras une note historique de fond en disant ce fond a été produit par les occupants allemands, de telle année en telle année.  En archivistique, tu as forcément ça. Et nous, on a donc on a rajouté une note, peut-être qu'il faut qu'on le fasse pour Radio Paris aussi, qui est pareil, qui est la radio de collaboration. Et il peut y avoir quelqu'un qui tombe sur un document Radio Paris. Ça peut être des fois un stagiaire aussi dans une boîte de prod.  Et il ne va pas se dire que Radio Paris, c'est la radio de collaboration, tu vois. Et il peut très bien se dire, là, ils sont en train de dire ça : « Les Américains ont trop bombardé, c'est un scandale ».  Mais c'est pour ça que là, on sait qu'il y a une réflexion là-dessus.  On a mis des notes sur l’actualité française qui est visible sur ina mediapro. C'est ça la différence aussi avec l’archivistique, c'est que on n'a pas forcément cette tradition de notices.  Alors, on a des notices d’ensemble maintenant, dans totem, mais qui ne sont pas visibles à l'extérieur. Mais toi, tu dois connaître, par l’archivistique, qu'il y a cette tradition de faire des notices, qu’on appelle notices de fond, et historiques du fond. Et ça, c'est quelque chose d'important, quand on te donne un fond privé, on te dit en quelle année il a été arrêté, qui a filmé, quel est l'organisme qui a payé en gros. 

\textbf{MF :} C'est vrai que même dans les archives papier, il y a toute une part importante qui est donnée à la remise en contexte. Parce qu'une archive sans remise en contexte... 

\textbf{PL :} Oui, c'est ça. C'est la base.  C'est pour ça que nous, on se bat beaucoup. Parce que même quelqu'un qui fait un sujet d'histoire, s'il utilise l'archive qui n'est pas la bonne ou dans un mauvais contexte, ça ne va pas parfois. Mais ça, c'est hyper important et les étudiants, je trouve que je leur fais réfléchir à ça aussi, toujours te poser la question, qui a filmé, qui a produit, toujours te poser des questions-là, même quand tu trouves les archives sur ina mediapro. Il faut toujours te poser ces questions-là, ce sont des questions de base, mais il faut toujours se poser.  

\textbf{MF :} Oui, donc il y a un peu une question de responsabilité de l'INA dans la contextualisation de ses archives. 

\textbf{PL :} Oui, on a aussi le fonds Marshall, c'est le ministère des Affaires étrangères. C'est le fonds que les Américains ont produit, c'était une unité américaine, c'est le fonds qui montrait comment les Américains étaient si parfaits après la guerre au niveau industrie, agriculture.  Donc c'est un fonds de propagande aussi, tu vois.  Et quant à ça, il faut toujours avoir cette idée que ce fonds, c'est un fonds produit par les Etats-Unis. Et sur ina mediapro, on a beaucoup de questions partenaires. Donc ça aussi, se poser une question partenaire, parfois, c'est filmé par MSF. Il ne faut pas oublier ça.  Mais en tout cas, aux jeunes générations, c'est important qu'ils se posent cette question-là aussi toujours. 

\section*{Les outils et services de l’INA }

\subsection*{Pascal Lebeurier}

\textbf{MF :} Les professionnels font leur demande sur INA Mediapro ? 

\textbf{PL :} Ça dépend, pas forcément. Soit il y en a qui, en fait, par exemple, pour le SAC, c'est plutôt une demande par mail, ils ont des portefeuilles. C'est-à-dire que quelqu'un s'occupe de la Cinémathèque, quelqu'un s'occupe de la mairie de Paris, quelqu'un s'occupe des grands musées. Là on travaille sur Orsay, qui va fêter ses 40 ans, c'est une demande intéressante. On a reçu une demande du Musée d'Orsay qui a dit « on va fêter nos 40 ans, qu'est-ce que vous avez de bien ? ». Et ensuite on fait une demande plus précise, mais c'est hors Ina mediapro.  Par contre après on leur renvoie le contenu sur Ina mediapro, pour qu'ils puissent visionner. Par contre sur les demandes commerciales, parfois le discours à l'INA dit aux clients commerciaux, quelqu'un qui prépare un documentaire ou une boîte de prod leur dire débrouillez-vous sur ina mediapro, autrement la recherche est payante. S'ils passent par un doc ils doivent payer la prestation.  Donc des fois ils se débrouillent tout seuls. Par contre toutes les productions INA, il y en a quand même beaucoup par an. Là on travaillait sur l'histoire du Téléthon par exemple. Là on a fait un truc sur le football français. Là, dans le contrat entre l'INA et le coproducteur, il y a une partie de recherche documentaire, qui a un nombre de jours, dix jours, douze jours.  Et donc là, il y a un documentaliste INA qui les accompagne. On a fait une grosse production cette année, qui était Renault, le chanteur Renault, qui dure 2h20. Et là, il y a une documentaliste qui a fait vraiment toutes les recherches, qui a envoyé vers le contenu au réalisateur. Et lui, après, il a tout dérushé. Donc il y a différentes pratiques mais parfois oui en effet ils peuvent se débrouiller sur une ina mediapro pour tout ce qui est commercial en effet. 

\textbf{ML :} Ça arrive des fois de ré-indexer quand vous avez recherché ?  

\textbf{PL :} Oui, on le fait quand on trouve et si il manque quelqu'un, on va l'ajouter.  Si il y a quelqu'un qu'on trouve, on va l'ajouter, on va corriger la notice.  Et puis les clients nous demandent eux-mêmes sur Ina médiapro. Sur Ina mediapro, tu as une zone « signaler une erreur ».  Il y a beaucoup de clients en Ina mediapro, des spécialistes de cinéma, qui nous envoient des corrections. Ils nous disent : « ah là vous avez oublié d'indiquer qu'il y avait Truffaut ». C'est contributif et on a plusieurs personnes qui envoient beaucoup de corrections et donc il y a des documentalistes qui les prennent en compte ça s'appelle la hotline documentaire et ils corrigent 

\textbf{MF :} Sur Ina mediapro, finalement, les professionnels n'ont pas, par exemple, les transcriptions. 

\textbf{PL :} Non, mais ça viendra peut-être dans le nouvel outil. Parce qu'en fait, ce qui est prévu, c'est qu'il y aura un nouvel Ina mediapro qui va venir. Nous, on va avoir Notilus, il y aura une nouvelle plateforme. Et il y a une des questions, est-ce qu'ils auront la sémantique ou la transcription ? C'est quelque chose qu'on envisageait. Donc il faudrait que tu poses la question peut-être à Xavier Lemarchand ou quelqu'un qui s'occupe des outils, c'est le spécialiste IA.  C’est quelque chose qu'on débat. Parce que peut-être qu'on ne mettra pas la sémantique dès le départ, c'est le débat, je crois que ce n'est pas tranché encore. Peut-être qu'on ne mettra que la transcription, pas la sémantique. Mais la transcription, sans doute on la mettra. Parce que ça apporte un plus, et ça marche bien. La sémantique, je ne sais pas.  Peut-être que ça restera au début un petit INA. Mais ça, ce n'est pas tranché encore. 

\subsection*{Arthur Hénaff} 

\textbf{AH :} Le service INAthèque est chargé de mettre en consultation l'intégralité des collections de l'INA pour les chercheurs, c'est-à-dire à la fois les archives pro, le dépôt légal Radio / TV  / WEB, les archives écrites et la documentation écrite. Et donc c'est pour les chercheurs au sens du dépôt légal, c'est-à-dire un sens très large, ça peut être des recherches académiques ou des recherches personnelles ou des recherches professionnelles.  Donc on met ces collections en consultation et par ailleurs on accompagne les chercheurs et les différents prescripteurs dans la découverte, l'utilisation, l'analyse et l'exploitation de ces collections. On forme aussi des étudiants, des groupes d'étudiants qui nous sont amenés par leurs enseignants. À partir de la L3 principalement et du master. On forme ces étudiants à la fois à l'histoire des collections, à la nature des collections de l'INA, et puis à la manière dont elles peuvent être utilisées dans le cadre d'une recherche académique. On travaille par ailleurs à valoriser les travaux qui sont produits à partir des collections de l'INA, notamment en tenant un carnet de recherche « hypothèse », qui s'appelle le carnet de recherche de l’INAthèque, je crois, tout simplement, sur lequel on publie plusieurs choses. Les documentalistes qui écrivent des billets sur l'arrivée de nouveaux fonds, par exemple, et notamment de nouveaux fonds écrits, ou qui vont faire des billets sur la manière dont tel ou tel fonds a été traité, avant d'être mis en base de consultation. C’est pour montrer le travail archivistique en train de se faire. Il y a des chercheurs et des chercheuses qu'on accompagne ou bien des chercheurs extérieurs qui travaillent sur des sujets liés aux médias, qui nous proposent la rédaction de billets d'analyse en profondeur d'une thématique. Pareil, c'est dans l'idée de tester des hypothèses de recherche qui leur serviront après dans leurs thèses ou dans d'autres articles, dans des revues académiques. Par exemple, cette année, on a publié une suite de trois articles sur le traitement médiatique du cyclone Chido. On a des articles qui sont des retours d'étudiants et d'enseignants sur les dispositifs pédagogiques qu'on met en place pour eux.  Et puis des articles de découverte des collections de manière globale. Côté Lab, c'est une équipe avec des ingénieurs de recherche qui travaillent au Lab. Ce n'est pas qu'ils font exactement la même chose mais leur travail, c'est aussi la mise à disposition et l'accompagnement des chercheurs, non plus tellement pour l'exploitation et l'analyse des médias en eux-mêmes et des notices documentaires rattachées à un média particulier, mais pour l'exploitation des jeux de données de l'INA, donc sur des grands datasets, sur des corpus beaucoup plus importants.  Le Lab accompagne ces chercheurs à la fois pour la mise à disposition de ces jeux de données et propose un accompagnement théorique et pratique sur la manière dont on fait des humanités numériques à partir de données audiovisuelles. Il les accompagne aussi sur l'utilisation des outils qui peuvent leur être utiles. Et le Lab a un séminaire tous les deux mois, où on reçoit notamment des gens qui sont ou qui ont été accompagnés par le Lab et qui nous font un point d'étape sur leur recherche.

\textbf{MF :} Du coup, dans l'objectif de ce point d'étape, ils peuvent récupérer d'autres informations à ce moment-là ?  

\textbf{AH :} Oui, l'idée, c'est que le séminaire, qui est ouvert à tout le monde, soit l'occasion de discuter.  Et donc, souvent avec du public extérieur ou du public INA, on essaie d'engager une discussion, notamment sur des points méthodologiques, de suggérer des corpus auxquels ils n'auraient pas forcément pensé, de suggérer des méthodes d'analyse auxquelles ils n'ont pas eu encore recours, ou de questionner ce qu'ils présentent dans leur méthodologie. Parce que le but de ces présentations, c'est évidemment d'esquisser des premiers résultats, mais c'est aussi de mettre à l'épreuve de la discussion ces méthodologies. Et chaque séance, il y a deux présentations de ce type-là. Chaque séance commence par une présentation qui est un peu plus une forme de « Keynote », où on fait venir un chercheur ou une chercheuse confirmée qui vient présenter des aspects plus conceptuels du sujet retenu. L'année dernière on a eu par exemple des séances sur sur la santé, sur le climat, sur l’emploi, sur les mobilisations sociales, et sur l'interventionnisme politique dans les médias.

\textbf{MF :} Est-ce que ça arrive qu'à cette occasion, vous voyez qu'il y a certaines erreurs ou certaines incompréhensions dans les données et que des documentalistes y reviennent derrière ? 

\textbf{AH :} Alors, en fait, côté lab, il y a deux choses. Il y a les ingénieurs de recherche qui, eux, vont vraiment se positionner sur la manière dont on traite les données.  Et il y a un accompagnement documentaire qui peut être fait pour justement comprendre comment ces données ont été construites au départ, ce qui relève des imports, ce qui relève des métadonnées qui sont produites par les documentalistes eux-mêmes, etc. Donc il peut tout à fait y avoir des discussions sur ces deux aspects-là.  Après, ce n'est pas tellement dans le cadre du séminaire qu'on va avoir une discussion sur des erreurs manifestes, parce que justement le travail est quand même fait en amont. Les chercheurs qui ont besoin de travailler sur des données INA, qui ont besoin d'un accompagnement, côté lab sont reçus, il y a des permanences toutes les semaines, donc ils peuvent tout à fait venir solliciter les collègues du Lab qui justement, ont fait leurs propres enquêtes documentaires et leurs enquêtes d'ingénierie, on va dire, et qui mettent les chercheurs sur la bonne route.

\textbf{MF :} Est-ce qu'on peut revenir sur les outils qui sont à la disposition des chercheurs 

\textbf{AH :} Il y a plusieurs choses. Côté INAthèque, ce sont des outils de recherche et de consultation.  On fait une distinction à l’INAthèque entre ce qu'on appelle les outils experts et les outils autonomes.  Juste pour rappel, l’INAthèque, c'est une sorte de réseau de consultation, avec une salle tête de pont, une salle principale, qui est la salle au sein de la BNF, qu'on partage avec le service son vidéo multimédia de la BnF, et où les chercheurs sont accompagnés et accueillis par une équipe de documentalistes.  Et là, il y a des outils experts, de même que dans les délégations, puisque dans les six délégations INA, il y a aussi des responsables documentaires qui peuvent faire un accompagnement particulier. On a l'outil de recherche en lui-même qui s'appelle « Hyperbase », qui permet de chercher dans toutes les bases de consultation des collections, qui permet de faire des paniers de documents audiovisuels qui sont ensuite consultables dans un autre outil qui s'appelle « Mediascope », qui permet de faire une forme de segmentation et d'annotation des médias. Par ailleurs, à partir de la liste des résultats, les chercheurs peuvent faire des corpus à partir d'un outil qui s'appelle « Mediacorpus ».  Ce sont des outils très experts qui ont été pensés à l'origine avec les chercheurs et en fonction des besoins des chercheurs.  Mais ce sont des outils aujourd'hui qui sont très anciens, qui ont plus de 20 ans, qui tournent sous les OS qui sont assez anciens. Et de notre côté, c'est vrai qu'il y a un besoin de renouveler ces outils. Et par ailleurs, l'outil autonome, c'est ce qu'on a déployé dans nos postes de consultation multimédia. L'INA a fait le choix, qui à mon avis est un très bon choix, de déconcentrer, de décentraliser la diffusion de son dépôt légal et de conventionner avec une cinquantaine de lieux partenaires qui sont des médiathèques, des bibliothèques universitaires, des cinémathèques et des archives départementales, de conventionner avec ces lieux pour mettre à disposition des postes sur lesquels des particuliers peuvent venir faire des recherches et consulter.  Là-dessus, ils ont accès à la fois à une interface qui s'appelle le « WebMedia », qui leur donne accès au dépôt légal du web, et à une autre plateforme, peu importe, c'est « OpsisMedia » en V1, et on est en train de déployer la V2 avec « Skopus », pour tout le reste, donc dépôt légal, radio, télévision et sources écrites. Et ça, le PCM V2, on va dire, cette deuxième version de l'interface radio-télévision CNC et archive écrite, on est en train de finir de la déployer cette année et on essaye d'y intégrer, de manière expérimentale, des fonctionnalités qui correspondent plus aux pratiques des chercheurs aujourd'hui donc j'anticipe un peu sur une question qui va venir plus tard mais typiquement aujourd'hui dans tous les outils de consultation de collection en bibliothèque, par exemple, il y a un besoin d'avoir un compte personnel, ce qui n'existait pas avant avec les PCM.  Un compte personnel sur lequel on puisse sauvegarder ses recherches, sur lequel on puisse faire des dossiers, et à partir duquel on puisse faire des corpus, corpus à partir desquels on puisse commencer à faire quelques analyses à très grands traits, quelques data visualisations.  Donc ça, ce sont des choses qu'on a inclus dans cet outil-là. Et de toute façon, il y aura un objectif, quelle que soit la trajectoire que prendra la refonte de nos outils, d'intégrer des fonctionnalités qui correspondent aux besoins des chercheurs aujourd'hui.

\section*{L’expertise du documentaliste }

\subsection*{Documentaliste} 

\textbf{Doc :} Quand on reçoit les étudiants à l’Inathèque, c'est notre rôle de les alerter sur le fait que, selon les périodes de diffusion, selon l'ancienneté, il peut y avoir des incidences. Le résultat de recherche peut être différent selon la façon dont on pose la question. Donc, choix du vocabulaire, évidemment, même si, quand on fait de la reprise d’antériorité, on utilise un vocabulaire sans doute plus récent et pas le vocabulaire de l'époque.  Et l'évolution des outils, parce que quand on avait les petites pages cartonnées...  Elles sont dans la base de données, mais elles n'ont pas été toutes complétées. Par exemple, si la notice n'a pas été complétée, il n'y a pas de mot-clé. Donc on sait qu'avant 1975, il faut parfois faire une recherche tout texte et pas uniquement par index mot-clé. Si on ne cherche que par index mot-clé, on risque de ne pas faire ressortir un résultat. En radio, pareil. Il y a beaucoup de notices où il n'y a qu'un titre.  Donc si on fait une recherche par mots-clés, on passera à côté de certaines notices. C'est un peu évident, parce que Google, c'est du tout-texte aussi. On sait que parfois, si on ne veut pas avoir trop de résultats, on peut cibler par les mots-clés. En se disant, par mots-clés, on ne va ressortir que les notices vraiment pertinentes qui parlent de notre sujet. Mais ça, ça fonctionne sur les fonds plus récents. On sait que sur les fonds anciens, il faut parfois lancer un grand filet pour ramasser quelque chose, peut-être parfois trop, mais en tout cas, pour être sûr de ramasser des poissons. Donc selon les périodes qu'on interroge, on va avoir une stratégie un peu différente.  Et ça, ce sont des choses qu'on ne peut pas deviner. Ce sont des choses qu'on doit transmettre, justement.  Alors ça, c'est pour les fonds anciens, mais par exemple pour les fonds plus récents, quand on parlait de l'évolution de l'indexation, par exemple, l'arrivée d'imports. Les imports d’une société prestataire qu'on appelle aujourd'hui « Onclusive », mais qui s'appelait Kantar, ou encore avant TNS. A partir de 2008, cela nous a permis de mettre du contenu sur les chaînes d'information continue sur lesquelles on n'avait pas les moyens de travailler, avant la transcription automatique. Mais on sait que ce n’est pas forcément fiable à 100\% et que ça fait beaucoup de bruit. Donc quand on fait une recherche sur ces chaînes-là, on sait qu'on va avoir plein de bruit et qu'on va être un peu noyé mais que c'est normal. Et ça, on doit pouvoir expliquer aux étudiants pourquoi ils se retrouvent avec beaucoup de notices un peu répétitives. C'est bizarre mais c’est ça ou rien. Alors, évidemment, on ne raconte pas tout aux étudiants la première fois. C'est en fonction de leur sujet de recherche.  Selon […] la période sur laquelle ils travaillent, ou si c'est une période transversale, on leur donne quand même quelques billes pour éviter de faire fausse route. Parce qu'on peut très bien interroger la base, avoir quelques résultats, et se dire, c'est bon, OK. Alors que ce sera peut-être faux. Ou partiellement erroné, parce qu'on n'aura peut-être pas posé les bonnes questions de la bonne façon. Donc, en plus de l'histoire de la radio, de la télé, il faut assimiler l'histoire de l'archivage à l'INA, l'évolution des supports, l'évolution des logiciels. Ce sont des choses qu'on apprend sur le tas, en fait. Parce qu'il n'y a pas de manuel pour ça. Moi, j'apprends toujours en interrogeant la base : parfois je me demande pourquoi je trouve, pourquoi il n'y a rien qui sort, il y a un truc qui ne va pas. Voilà, on découvre toujours, on apprend toujours des choses comme ça.

\textbf{MF :} Donc il y a tout un travail pour comprendre les bases de données, mais du coup, est-ce qu'avec le temps, ce n'est pas une connaissance qui risque de se perdre ? 

\textbf{Doc :} Ça, c'est le risque aussi. C'est notre rôle de le transmettre. On sait que nous-mêmes, on nous a appris des choses. Les anciens, les anciennes, nous ont transmis des choses et maintenant, c'est à notre tour. Mais on en a certainement perdu en route.  Il y a des connaissances sur la façon d'indexer à l'époque, ou des choses qu'on a peut-être transformées dans notre mémoire. Je pense que ça, c'est inhérent à l'histoire de toutes les archives.  Quand on travaille en histoire, sur l'époque médiévale, je suis sûre qu'il y a des questions semblables. On voit un truc, on se demande pourquoi c'est comme ça, mais on ne sait pas.  A l’INA, c'est pareil, à un moment donné, on se dit, mais pourquoi c'est indexé comme ça, pourquoi là, cette collection, elle est indexée sur cette telle période, et puis là, il n'y a rien, et parfois, on ne sait plus pourquoi.

\textbf{MF :} A propos de l’évolution du métier de documentaliste :

\textbf{Doc :} Là, pour le coup, le fond du fond, ça n'a pas tellement évolué. Il faut quand même être assez rigoureux, être curieux parce que quand tu fais de l'indexation, tu traites de tous les thèmes d'actualité. Donc il faut un peu de culture générale pour comprendre de quoi on parle. Avec un peu d'histoire, je pense que ça fait partie des éléments habituels, ça, ça ne change pas. Même pour répondre aux étudiants. Ils travaillent sur toutes sortes de sujets. Même si on ne connaît pas tout, évidemment. Quand les gens s'inscrivent, souvent on a leur sujet. Généralement, on prépare un peu à l’avance. On regarde un peu sur internet de quoi il retourne, si on ne connaît pas, ça arrive. Donc, on se renseigne un petit peu avant pour pouvoir ensuite les guider dans nos bases. Ça fait partie du travail en amont. Qu'est-ce qu'il y a à dire sinon par rapport au métier ? La formation, elle a pas mal évolué, parce que comme je te dis, moi, quand je suis rentrée à l'INA, il n'y avait pas de formation encore de documentaliste audiovisuel. Les gens apprenaient sur le tas, et puis bon, au sein de l'INA, c'est l'INA qui formait, parce que l'INA avait déjà cette mission de formation depuis l'origine, donc elle formait des techniciens, mais aussi des documentalistes en interne, mais uniquement à ses outils, évidemment. A la création de la première formation, je ne sais plus, je crois en 1999, je ne sais plus trop... à partir de ce moment-là, l’INA n’organise plus de concours.  L'accès se fait par un recrutement classique. Cette formation est au départ en partenariat avec le CNAM. Ensuite, ils se sont séparés. Le CNAM a continué à faire sa formation. Et l'INA a créé INA SUP, avec une licence professionnelle 

\subsection*{Pascal Lebeurier }

\textbf{ML :} Oui, donc c'est quand même encore une part d'expérience importante. 

\textbf{PL :} Oui, c'est ça, il y a ça aussi qui est important (montre les fiches physiques). 

\textbf{ML :} Vous recherchez encore dans ces casiers ? 

\textbf{PL :} En fait, il y a des documentalistes très expérimentés, ils ont travaillé à la prod, ils peuvent très bien parfois aller jeter un œil au fichier. Il y a Malika, Nicolas Coste, qui est un documentaliste qui a beaucoup d'expérience, ont encore des fois le réflexe d'aller au fichier.  Voilà, donc c'est bien qu'ils soient encore là. On a des documents d'accompagnement aussi papier, parfois, des classeurs, là où il y a les conducteurs. Et il y a les documents d'accompagnement, c'est-à-dire les rapports en chef de chaîne, ça, on jette [un coup d’œil] parfois.  Quand on se dit, tiens, c'est bizarre qu'il n'y ait pas eu d'émission avec telle personne, tel jour, on va regarder la prochaine chaîne, dans les vidéoprogrammes, dans les Télérama. Donc, par exemple, pour le ministère d'archives avec Serge Gialet, je le faisais souvent, on allait chercher dans les documents d'accompagnement. 

\textbf{MF :} Donc il y a quand même plusieurs sources d'informations pour un seul et un même document. 

\textbf{PL :} Oui, c'est un peu ça. Et puis Notilus, ça nous donne des transcriptions, c'est intéressant.  Et puis il y a les sémantiques qu'on commence à utiliser, mais là c'est un peu différent. La recherche sémantique, ça va te donner des idées, il va te proposer des choses auxquelles tu n'aurais pas pensé, il va te donner des idées. Tu te dis « ah oui, tiens, je pourrais chercher tel mot », et après tu retournes parfois dans Totem pour chercher que le mot. C'est un outil. Le fichier papier c'est quelque chose qui reste, malgré tout, moi j'en ai fait quand j'étais documentaliste, on a passé du temps à remplir dans Totem tout ce qui était en fichier papier.  Mais par contre à l'époque on ne visionnait pas. Donc ça veut dire que parfois on a rempli des choses qui étaient là. Finalement quand tu visionnes l'extrait, tu ne retrouves pas exactement la même chose. Mais bon, tout ça c'est intéressant. Ça a été fait par les cinémathécaires de l'ORTF. D'ailleurs, il n'y a pas longtemps qu'on en a rencontré une, quelqu'un à l’INA ou sa belle-mère, était une des cinémathécaires de l’ORTF. C'est toujours intéressant de dire qu'elles ont travaillé à l'époque. Il y a beaucoup de fautes de frappe, par exemple, en foot, il y a des énormes fautes de nom, mais c'est normal, à l'époque, ils avaient très peu d'infos. 

\textbf{MF :} Oui, puis ça devait aller un peu plus rapidement, parce que c'était en temps réel.   

\textbf{PL :} Oui, oui. Notilus c'est intéressant, mais les fichiers papiers ce sont aussi des outils qui restent très intéressants. 

\textbf{MF :} Même en ayant tous ces fichiers qui sont numérisés ? 

\textbf{PL :} Oui mais tu peux interroger Nicolas Coste, qui est un des documentalistes, il aime bien aller vérifier dans le fichier papier, parce qu'il est habitué, il travaillait à la prod. Là, par exemple, on travaille sur le club de basket de Pau-Ortez. Je sais que, c’est les années 70, quand tu n'as pas un stade 2, il faut aller vérifier. Parfois, le matériel est mal rattaché. Le reportage peut exister, mais la notice est sans matériel, finalement, le matériel existe ailleurs. Il y a toutes les vérifications à faire. Tu peux aller vérifier dans les rétro-régionales. L'équipe de Pau, c'est le Bordeaux, c'est le fonds Aquitaine. Parfois, tous les matériels sont simplement un titre.  Il y avait un hebdo, hebdo actualité, hebdo sport Aquitaine, qui était une fois par semaine et il n'y a rien, il n'y a pas de descriptif.  Donc tu sais que dans l'hebdo, tu auras peut-être les images du basket de Pau. Il y a tout un truc à aller faire. Et ça, si tu cherches la transcription, avec le Pau-Ortez basket, je ne suis pas sûr que les résultats soient forcément très probants. Il y aura moins de bruit, ça dépend comment le journaliste le dit. En tout cas, il faut le savoir que tu as des notices mal indexées. Notamment aussi en radio. Le fond radio est encore moins décrit que le fond télé. Pour la transcription, pour la radio, par contre, ça marche très bien. 

\textbf{MF :} ça peut vous arriver que des professionnels reviennent vers vous parce qu'ils n'ont pas trouvé certaines informations ? Par exemple, avec des fiches anciennes où il n'y aurait pas d'indexation ? 

\textbf{PL :} Oui, ça arrive qu'ils reviennent vers nous. Parfois, ils ont du mal. Les réalisateurs, quand ils cherchent eux-mêmes, ils le disent eux-mêmes, ils ne sont pas de grandes documentalistes.  Donc ça veut dire que tu trouves des choses en tapant le nom celui-là ina mediapro. Mais il y a certaines choses, il faut qu'il passe par un documentaliste. Et les réalisateurs, ils préfèrent dès le départ. Tu vois celui sur Renault, Tancrède Ramonet, qui est un très bon réalisateur. Il a dit, moi, je veux passer par un doc. Et il a demandé à un documentaliste, tu me sors tout ce qu'il y a sur Renault dans les années 70. Moi, je veux tout voir, je ne veux pas passer à côté d'une archive intéressante. Il avait besoin de tout visionner. J'ai un autre réalisateur, pareil. Il a travaillé sur Raymond Devos, on a fait un docu pour la première, il avait demandé un doc.  Parce qu'ils ont raison, ils disent qu'eux, ils sont réalisateurs, ils ne sont pas documentalistes.  Un documentaliste, il va trouver des choses que eux ne trouveront pas. Parce que le doc, il sait chercher autrement, il connaît les émissions, il connaît les fonds, il sait quelles émissions vont être bien, sur lesquelles il ne faut pas passer à côté. Donc, ça l'expérience, on appelle ça l'expertise du documentaliste, elle est hyper importante. La nouvelle chaîne de vente, d'ailleurs, il est prévu, peut-être que les grandes forces qui disent aux clients, vous pouvez trouver vous-même, autrement vous passez par les services d’un doc. Donc ça vous coûtera plus, mais vous trouverez des choses. Tu auras peut-être le truc vraiment plus facile, la haute culture, si on veut, mais c'est prévu qu'il y aura peut-être ces distinctions.  

\textbf{MF :} Mais du coup, est-ce qu'avec le temps, il n'y a pas un risque que ce savoir et cette expertise se perdent un peu, si ça se renouvelle ? 

\textbf{PL :} Non, non, je ne pense pas, au contraire. Parce qu'en fait, il y a des documentalistes qui travaillent beaucoup maintenant aussi sur la connaissance des fonds. Après, alors peut-être que le risque, c'est que ça arrive, par exemple, pour la musique. On a des docs musiques à l’INA, parce que la musique classique ou le jazz, c'est assez spécifique. C'est Brigitte Decroix qui a une expérience énorme de la musique, elle va partir sûrement dans un ou deux ans. Elle forme actuellement des nouveaux docs, mais jamais les nouvelles personnes auront sa connaissance des fonds musicaux parce qu'elle est là depuis 40 ans, 30 ans. Donc il y a le risque que des personnes en effet partent à la retraite, mais toute personne qui part à la retraite, part avec une connaissances que peut-être quelqu'un d'autre n'aura pas forcément. Alors après, on capitalise. Une personne a travaillé sur les fictions, on a un corpus fiction maintenant. Dans Totem, tu as des campus émissions. On a des émissions pour la jeunesse, fiction, et la personne qui a fait les fictions, elle partie à la retraite, mais avant de partir, elle a fait un énorme travail sur les fictions : les téléfilms, les feuilletons. Donc ça, c'est capitalisé dans Totem. Il y en a qui connaissent bien le sport.  Moi, quand j'ai travaillé, beaucoup sur le sport, j'ai quand même fait de la connaissance des fonds, j'ai fait des corpus. Mais c'est comme dans toutes les entreprises.  Avant de partir, Brigitte, c'est ce qu’elle me dit en musique, son grand truc, c'est comment elle va donner le maximum de connaissances avant de partir. Donc c'est pour ça qu'on l'a fait travailler avec d'autres documentalistes. C'est bien parce qu'elle leur donne des connaissances, même des process, mais ils n'auront jamais sa connaissance. Elle fait de la musicologie. Mais bon, c'est comme ça. 

\textbf{MF :} Ce qui ressort de tout ça, c'est que les bases de données sont plus un outil où on va récupérer les données, récupérer les images, et que niveau informations sur ces documents, c'est plus les documentalistes qui les ont.   

\textbf{PL :} C'est ça, quand tu recherches à Gaumont-Pathé ou à l'ECPAD, il vaut mieux parfois demander aux documentalistes. Moi, j'en ai fait déjà.  Nous, on peut faire des recherches en ligne sur leur base. Il suffit d'avoir un compte Gaumont, par exemple. Tu fais des recherches, mais tu t'aperçois que c'est un fonds clos. Et si tu demandes à un documentaliste de Gaumont de faire la recherche, il connaît parfaitement son fonds.  Alors que toi, qui est doc INA, tu as été moins habitué, c'est sûr qu'il y a des spécialistes. Mais parfois, on recherche dans d'autres bases, pour certains, quand on travaille pour des projets. Parfois, on demande de chercher dans d'autres plateformes. Ça peut être aussi le musée de l’Air, ça peut être les archives du NARA, aux Etats-Unis, le NARA, c'est les archives américaines, nationales américaines. Il y a une base en ligne avec des photos, des vidéos. Mais on sait très bien qu'on cherche en ligne. Ce n'est pas choses simples. Bon c’est en anglais, pour ça on s'en sort. Mais en fait, on sait très bien que la meilleure façon de trouver des choses au NARA, c'est d'aller là-bas. Serge Viallet, par exemple, a trouvé des choses sur le débarquement, que même les Américains ne [connaissaient pas]. Parce qu'ils ont encore des supports physiques là-bas qui ne sont pas décrits. 

\section*{L’impact des nouvelles technologies – de l’IA}

\subsection*{Documentaliste} 

\textbf{Doc :} La transcription automatique, c'est bien pour certaines choses, mais je pense que le regard humain...  Par rapport à la mise en forme aussi de l'émission, je parlais d'un beau noir et blanc. Il y a des choses, on ne va pas dire dans la notice que c'est un beau noir et blanc, mais on va faire en sorte que ça puisse donner envie, de décrire un peu la forme de l'émission, parce que ça, l'IA ne sait pas encore très bien faire. Je pense que dans l'avenir, sans doute notre apport premier, ce seront tous ces éléments de connotation, tout ce qui peut transparaître dans l'ambiance d'un débat ou la construction d’une émission, par exemple si c'est un peu plus animé ou virulent, même si nous, dans notre description, on essaie d’être le plus objectif possible. Il y a des moments, s'il y a un incident, on le signale quand même, dans un champ à part. Dans « Aujourd'hui Madame », par exemple, j'ai entendu des propos qu'on qualifierait aujourd'hui de pédophiles. Quelqu'un dit, un intellectuel, qu'il considère que le sexe avec les enfants, oui, c'est bien. Et ça, au détour d'un débat, aujourd'hui, c'est « gloups ».  Du coup, je ne savais pas trop comment tourner la chose. Je l'ai signalé quand même dans le champ « Notes », en citant les propos de l’intervenant, et dans mon choix de mots-clés. Bon, à ce moment-là, il n'y avait pas encore la transcription automatique. En tout cas, je l'ai signalé.

\textbf{MF :} Oui, parce que ça peut être d'intérêt. Même si l'IA arrive et semble pouvoir changer beaucoup de choses dans le métier, il y a toujours l'intervention humaine qui est super importante. 

\textbf{Doc :} Et puis, parce que dans l'histoire de l'IA, il y a la question de la fiabilité. C'est que l'INA est censé être un référent. Donc, on sait bien que même si nous, dans notre description, notre travail, il y a forcément une part subjective qui est loin d'être parfaite, ça, c'est sûr : il y a des fautes d'orthographe, il y a une forme de subjectivité. Dans ce qu'on retient d'une émission quand on fait un résumé, il y a forcément une part de subjectivité, c'est inhérent. Je sais que d'un collègue à l'autre, on ne va pas décrire de la même façon. Donc l'IA, ça va être pareil, il y aura une marge d’erreur. Le problème de l'IA, c'est que si elle fait des hallucinations, si elle invente des choses qui ne sont pas, ça pose problème quand même. Parce que nous, on peut peut-être omettre quelque chose. De fait, comme je disais, l'histoire de la pédophilie et tout ça, à l'époque, ça n'avait peut-être pas été retenu. Déjà, il ne faut pas oublier que c'était indexé en direct à l'époque. Dans le flux de la conversation en plateau, ce n'était pas forcément retenu. Ce n'était pas forcément un thème non plus à l'époque. Du coup, il y avait des choses sur lesquelles on passait, dont on ne parlait pas, par omission. Mais on n'allait pas inventer quelque chose qui n'existait pas. C'est là-dessus que j'ai un peu peur. Pour faire un contrôle qualité, je ne sais pas comment on peut faire sans visionner. On peut imaginer peut-être des outils qui seraient capables de vérifier des incohérences.  Je pense que l'IA peut vérifier des incohérences, mais si c'est cohérent, comment elle peut détecter une erreur ? Enfin, si elle invente un truc qui n'est pas...  Par exemple, au début, je me souviens, il faut demander aux collègues qui ont travaillé sur l'émission de France 5, « C dans l’air », pour rédiger des prompts d’indexation automatique. Quand ils avaient fait leurs premiers tests, dans les résultats, on aurait pu penser que l’émission parlait de terrorisme tout le temps. Alors qu'en fait, pas du tout, c'est juste parce qu'il devait y avoir un mot qui faisait que la machine le reliait au terrorisme. Donc après, ils ont dû affiner le prompt et cela a fini par fonctionner. Mais ça suppose un travail dingue. Donc, s'il faut regarder l'émission et après relire le truc, enfin je veux dire, le temps passé, à un moment donné, je me dis « où est le temps gagné ? ».  Enfin, vraiment, je ne sais pas.

\textbf{MF :} Oui, il y a encore beaucoup d'élaborations nécessaires.   

\textbf{RD :} C'est ça. Malgré tout, l'intérêt pour moi, c'est quand même pour les choses qu'on ne traite pas. Comme les premiers imports en 2008, c'était pour les chaînes d'info continue sur lesquelles on n'avait pas le temps de passer. Ça sera pareil, parce que là, il y a plein de chaînes sur lesquelles on ne travaille plus du tout. Avant, on traitait Europe 1, RTL en radio, RMC. Bon, RMC, je n'avais pas une passion pour RMC, mais en même temps, ça permettait de savoir de quoi ils parlaient. On avait une connaissance des médias qu'on perd nous-mêmes. Parce que de fait, on ne passe plus dessus. Pareil pour les chaînes du Câble Satellite, LCI, BFM, France 24, on traitait toutes ces chaînes-là et on ne le fait plus. Moi, je ne connais plus leur grille, sauf si je regarde un peu chez moi. Là, l'IA aura un rôle pour compléter, parce que ça permettra quand même, plutôt que de n’avoir rien, ça permettra des recherches sur ces chaînes. Pour nous, c'était ingérable, on ne peut pas traiter un truc 24 heures sur 24. Il faut que ce soit complémentaire.   

\textbf{MF :} Oui, ce serait l'idéal.  Parce que du coup, si ces chaînes sont plus traitées, c'est pour quelles raisons ?  

\textbf{Doc :} Je pense que c'est une question de moyens, de personnel, on est de moins en moins nombreux à faire ce travail-là. Mais, à terme aussi, est-ce que la question économique… parce que payer moins de personnel, ok, mais faire travailler l'IA, ça coûte cher aussi.  Donc je ne sais pas économiquement quelles sont les solutions.

\textbf{MF :} On est dans une période de changement.   

\textbf{Doc :} Oui, je pense qu'on ne sait pas trop. […] On parlait aussi du coût écologique, tout simplement environnemental. Parce qu'il faut avoir des grosses machines, des gros d'entrepôts, machin, qui consomment de l'énergie, qui va payer cette consommation-là aussi ?

\textbf{MF :} Oui, et puis même si c'est délocalisé, il y a toute la question de la gouvernance de la donnée, à qui ça appartient vraiment, de qui on dépend. 

\textbf{Doc :} Et puis après, le regard dessus. Parce que même si on dit que ça a été généré par IA, on sait que ça ne sera pas très fiable. Au bout d'un moment, il faudrait voir le retour des étudiants et des chercheurs. Qu'est-ce qu'on fait avec ça aussi ? Parce qu'au bout d'un moment, si on se casse le nez à chaque fois qu'on a le résultat et que ce n'est pas pertinent, à quoi ça sert ? C'est ça la question. Il y a des limites.

\textbf{MF :} Double dépense au final.   

\textbf{Doc :} Oui, c'est ça. 

\textbf{MF :} Actuellement, quels outils ou peu importe quels moyens pour toi améliorerait ton quotidien, ton travail ? Avec les données, si tu en vois. 

\textbf{Doc :} Dans ma pratique, […]. Personnellement, je n'ai jamais pris le temps d'utiliser Copilot.  C'est un peu idiot parce que bon, je pourrais. En fait, moi, tu vois, j'aime bien écrire. J'aime bien rédiger, j'aime bien résumer et du coup, je n'ai pas envie, j'aime bien travailler avec mon cerveau. Et le fait de faire une requête, ça me fatigue en fait. Je préfère réfléchir moi-même et puis pondre moi-même mon truc plutôt que de générer un truc qu'il faudra reprendre et qui n’est pas de moi, ça me gêne. Après, je comprends qu'on puisse l'utiliser pour certains trucs pour lesquels on n'a pas le temps. Tu vois, quand tu veux faire des accroches aux commandes, pour les réseaux sociaux, je peux comprendre qu'ils ont un temps très limité. Mais moi, personnellement, dans mes activités, en tout cas, je travaille sur des fonds, les inventaires, les dépôts de fonds anciens. Alors pour ça, l'IA ne peut pas faire grand-chose. […] Puis le flux, je connais. Ce sont des vieux outils, donc pour l'instant, je n'ai pas d'usage. Et dans ma vie personnelle, non plus. Après, on y est tous confrontés. J'utilise le moteur de recherche Qwant, où il y a un petit module d'IA automatique. C'est une réponse qui fait une synthèse de plusieurs sources. L'avantage de Qwant, c'est qu'il cite les sources, tu sais d'où ça vient. Ce qui est rigolo, c'est quand tu connais bien un sujet et que tu lui demandes un truc, tu te rends compte que ça peut dire n'importe quoi. Mais de temps en temps, c'est bien aussi. De temps en temps, ça marche quand même. En tout cas, ce genre de choses, ce sont des petits outils.  Mais pour l'instant, ça ne va pas plus loin que ça, je ne m'en sers pas vraiment. Après, voilà, c'est une question d'usage, d'habitude aussi.  Peut-être un peu de flemmardise aussi. Mais je ne suis pas contre, contre. Je vois bien l’intérêt, comme je t'ai dit, pour traiter des choses qu'on ne traite plus de toute façon. Après, si un jour ça mange complètement notre travail... pour l'instant, il y a une frontière qui est fragile.  Si, à terme, c'est ce qui se profile quand même, qu'il y ait des résumés automatiques et qu'ensuite on soit amené à les corriger, là je trouverai ça moins rigolo. Je serai peut-être à la retraite d'ici là, mais ça va peut-être venir plus vite qu'on pense. En tout cas, pour l'instant, je savoure les derniers moments de l'indexation avec des vrais morceaux d'indexation. Parce que moi, je dis toujours, même si ce sont des programmes pas intéressants, enfin, un peu chiants, nuls quoi, des fois, des choses bêtes. RMC, c'est insupportable à écouter, parce que c'est plein de pubs, ça hurle, c'est que du sport, tout ça, c'est gavant, mais en même temps, ça m'intéressait quand même de le traiter, parce que ça permet de, voilà, on voit ce que c'est. Sinon, chez moi, je ne l'écoute jamais. Et puis, c'est toujours amusant d'essayer de tourner les choses. Enfin, on essaie de rester objectif, mais en même temps, on se dit, voilà, j'essaie de tourner mon résumé pour faire en sorte que, bon, on rend compte du truc, quand même, de ce qu'il y a dedans. Même un programme nul que tu ne regarderais pas chez toi, finalement, à rédiger, moi, ça m'amuse un peu.

\textbf{MF :} Oui, c'est un métier enrichissant.   

\textbf{RD :} C'est ça. On ne fait pas de critique d’une émission, mais le fait d'essayer d’en rendre compte, déjà rien que ça, en étant prudent, rien que ça, c'est un travail du cerveau, une petite synthèse.  Et c'est quelque chose d'assez plaisant. Enfin moi j'aime bien, quoi. Après je sais qu'il y a des collègues qui ne supportent plus parce qu'au bout d'un moment ils trouvent ça trop répétitif et tout ça. Je pense que ça dépend de l'approche que l'on a, en fait. Il y a des moments, il y a des choses que je trouve chiantes aussi. Des fois, je suis un peu flemmarde, je me dis : j'ai pas très envie de faire ce truc-là. Oui, ça, c'est normal, on est tous comme ça.

\subsection*{Pascal Lebeurier }

\textbf{MF :} Pour l'instant, la transcription, c'est surtout sur les nouveaux documents ?  

\textbf{PL :} Non, ce sont les anciens. Puis on a des outils, on a Quicksand, si tu as besoin de transcrire. Et les recherches de la transcription, c’est dans les anciens aussi, on utilise ça. 

\textbf{MF :} L’IA a un fort impact actuellement ? 

\textbf{PL :} L'impact, c'est que le fait d'avoir transcrit une grosse partie du fond, ça nous aide en effet à trouver des documents qu'on n'aurait pas trouvés qu'avec Totem. Il n'y a pas longtemps, on nous a demandé le réchauffement climatique, depuis quand on en parle ? C’était pour « inattendu », une émission de France Info. La journaliste, elle est plutôt partie de 77, où il y avait un document qui s'appelait « Les mystères de la Terre ». On parlait déjà du réchauffement, mais quand tu tapes dans les transcriptions, tu t'aperçois que Paul-Emile Victor, dans les années 50, il parlait déjà du réchauffement quand il revenait du Groenland. 

\textbf{MF :} Mais du coup, dans l'indexation, c'est pas forcément [inscrit]… 

\textbf{PL :} L'indexation, ça va être marqué « expédition Groenland », il n'y aura pas le terme de réchauffement, alors qu'il emploie vraiment le mot réchauffement. Mais la transcription peut nous permettre de trouver ce genre de terme. Alors après, ça met du temps, il faut écouter. 

\textbf{MF :} Autant la politique documentaire a un impact sur l'utilisation des bases de données, autant il y a quand même une grande importance du savoir des documentalistes.  

\textbf{PL :} Oui, oui.  Après, nous, on sait très bien qu'on ne peut pas reprendre notre fond.  Il y a tellement de trous. Donc ça veut dire qu'il faut savoir qu'il y a des trous la transcription ça va nous aider beaucoup pour la radio notamment. Et pour les émissions qui arrivent, par exemple on va récupérer le mandat « C à vous », c'est pareil si on doit refaire 15 ans de « C à vous » en résumé, il faudrait des moyens humains. C'est pour ça qu'il y a aussi eu les représentations IA de résumé. Alors après, le problème de l'IA, c'est qu'un défi, il y a des erreurs, parfois.  Et puis, il y a toujours le problème : qu'est-ce qui fait perdre du temps ? Des fois, c'est plus la transcription.  On l'avait fait sur la gare d'Algérie, quand on avait fait la série pour Arte, on s'est aperçu qu'il y avait beaucoup d'erreurs, c'est les noms propres, les noms algériens, en effet.  Donc, on s'est aperçu qu'en plus, c’étaient les personnes qui avaient un accent.  Donc, ce n'était pas forcément efficace d'utiliser la transcription. Donc, les documentalistes, ont préférés plutôt écouter les entretiens. Et pour les entretiens patrimoniaux, ce serait intéressant aussi que tu interroges des docs qui ont indexé les grands entretiens patrimoniaux : « mémoire composée », « passé simple ». Il y a des analystes qui ont indexé, savoir est-ce qu'ils utilisent beaucoup la transcription ou pas. C'est que Madeleine en a fait, il y a beaucoup de docs. C'est l'équipe d'Emeline, je crois, qui le fait mais c'est intéressant de savoir est-ce que la transcription aide ou pas. 

\textbf{MF :} Oui, entre est-ce qu'en effet ça ajoute quelque chose ?  

\textbf{PL :} Si tu ne fais que vérifier, si tu passes énormément de temps à vérifier, parfois c'est compliqué. Il n'y a pas longtemps, on nous a demandé de vérifier la transcription.  On a une série sur la mode de la préhistoire, un docu. C'est un sujet très précis, la mode et la préhistoire.  Donc ils ont interrogé des paléontologues russes et allemands. Ils ont demandé de vérifier, parce qu'on a un doc qui parle russe ici, une traduction IA. La première fois, Charlotte, qui a fait l'allemand, m'a dit que ça a bien marché pour la première.  Par contre, pour la deuxième, il y avait trop d'erreurs dans la traduction IA. Donc il aurait fallu que je refasse dès le début. 

\textbf{MF :} Donc c'est prometteur mais encore à améliorer 

\textbf{PL :} Oui mais ça dépend des cas. Mais après, la transcription audio, à la radio elle marche très bien globalement ils ont fait beaucoup de progrès. La transition c'est un outil. Et les documentalistes qui l'utilisent pour l’indexation et traitement, ils disent parfois ça, que quand tu travailles pour une émission France Musique qui dure une heure, si tu te fais faire la transcription, ça va t'aider. Tu auras le nom des invités, tu auras le nom des œuvres. C'est vraiment l'outil qui est intéressant. Après, je sais qu'il y a des documentalistes qui l'utilisent, d'autres pas. C'est assez partagé. 

\subsection*{Arthur Hénaff}

\textbf{MF :} Dans ce cadre-là, c'est le métier des documentalistes qui évolue ? Pas forcément ceux qui conseillent à l’inathèque, mais ceux qui vont devoir revoir ces transcriptions. 

\textbf{AH :} Oui, tout à fait. Il y a des évolutions. Alors évidemment, je ne peux pas du tout me prononcer sur ce que vont être les évolutions du travail des documentalistes, mais des choses qu'on peut imaginer et qui en fait existent déjà. C'est que là où il y avait un travail qui se faisait beaucoup au moment de la production, dans le travail sur la qualité des données en amont au moment de la production. Aujourd'hui, il y a un travail qui va s'axer aussi sur le contrôle de la qualité de métadonnées documentaires ou de données de transcription produites par l'intelligence artificielle, par exemple.  Et par contre, à mon avis, le travail de documentaliste va aussi beaucoup, en tout cas à l’INAthèque c'est déjà le cas, et je pense de manière globale, s'axer sur toute la médiation qu'on peut faire autour de la connaissance des fonds et de la connaissance de la manière dont c'est traité, des avantages de tel ou tel type de traitement, notamment par IA et des limites aussi. Et ce sont ces explications-là, de contextualisation, qui vont être, à mon avis, même de plus en plus importantes. D’autant plus importante qu'on aura des masses de données elles aussi de plus en plus importantes. Parce que pour moi, il faut voir que trouver des résultats par une recherche plein texte via une transcription, c'est une première étape. Mais une transcription, elle ne renseigne sur rien d'autre qu'elle-même. C'est juste une transcription.  Et à mon avis, il y aura toujours à un moment ou à un autre un besoin de la mettre en regard d'une notice documentaire qui comportent des informations, à la fois historiques et techniques, qui vont permettre de remettre en contexte les documents qu'on a trouvés via cette recherche plein texte.

\textbf{MF :} Finalement c'est une évolution du métier, mais les notices en elles-mêmes restent très importantes pour leur mise en contexte ?  

\textbf{AH :} À mon avis, que ce soit sous la forme de notices ou sous d'autres formes, liées à la recherche sémantique, etc. Le fait de contextualiser une liste de résultats de recherche, c'est essentiel.  Après, ces éléments de contextualisation peuvent aussi eux-mêmes être assistés par une production IA. Par exemple, à partir de tous les épisodes d'une même collection, de faire des résumés de collection,  etc. Dans cette production-là, de différents niveaux de notices, de collections, de programmes, de transmissions et retransmissions, d’événements médiatiques. Je pense qu'il y a plein de choses qui sont testées, il y a plein de choses à imaginer, et à mon avis, le travail de médiation, justement sur tout ce qu'on fait et tout ce qu'on apprend des collections via ces traitements automatiques et tout ce qu'on apprend aussi grâce aux chercheurs et tout ce qu'on apprend par les traitements qu'on fait nous-mêmes et par le contrôle qualité qu'on fait nous-mêmes. C'est une médiation très humaine. Et à mon avis, il y a vraiment quelque chose à investir de plus en plus dans le métier des documentalistes là-dessus. Et en tout cas, c'est certain que c'est un objectif côté INAthèque de développer de plus en plus des profils qui soient en capacité de faire cette médiation auprès des étudiants et des chercheurs. Et notamment, de faire comprendre ce que c'est aujourd'hui et ce que ce sera demain que la description de contenu audiovisuel par IA. 

\textbf{MF :} Impact de l’IA sur les pratiques professionnelles ? 

\textbf{AH :} Alors, il y a évidemment toutes celles dont on a parlé et dont vous avez beaucoup plus parlé côté TEC sur le traitement documentaire. Mais nous, une partie importante de nos missions, c'est aussi la médiation auprès des étudiants. Et donc, c'est ce que ça veut dire d'utiliser des documents d'archives audiovisuelles dans un travail intellectuel, dans un travail de recherche.  Et aujourd'hui, on est vraiment confrontés à un défi qui est majeur qui est la production de fausses archives via les IA génératives. Ça, on en voit énormément sur les réseaux sociaux. On voit beaucoup de gens qui se mettent en scène dans des archives. On voit beaucoup de gens qui vont utiliser, par exemple, une photographie d'archives et puis qui vont l'animer. Et on voit beaucoup de gens qui produisent des fausses archives. C'est montré de manière plus ou moins assumée, on va dire. Parce que, même dans la production d'une fausse archive, le côté un peu « old school » de l'archive, toute cette patine patrimoniale que l'INA a vraiment rendue extrêmement sensible auprès d'un public, il y a une valeur patrimoniale qui est claire. Et des gens qui vont faire des fausses archives, font des fausses archives parce qu'ils savent que s'ils lui donnent un cachet patrimonial, un cachet ancien, un cachet un peu INA, ça va toucher plus de monde et ça fait un peu office de preuve, ça ancre dans l'histoire d'une technique, etc. Mais il y a vraiment cette question de la différence à faire maintenant entre les archives originales et les fausses archives. On va de plus en plus avoir besoin d'insister sur des points d'éthique de la recherche et de méthodologie de la recherche en l'audiovisuel avec ces IA génératives. Ce n'est pas parce qu'on a de l'IA que d'un côté, les transcriptions ou les métadonnées deviennent transparentes. Derrière, il faut une méthode de recherche, il faut des questions de recherche, il faut des hypothèses, il faut des corpus, il faut faire des choses dans le bon ordre. Et il faut faire aussi le même travail de sensibilisation sur un document audiovisuel. Nous, on doit aussi jouer notre rôle de tiers de confiance. Auprès de la population académique et du grand public qui utilise les collections à des fins de recherche. Donc il faut qu'on soit très attentifs à ce qu'on va devoir, à mon avis, repenser, réinventer en termes de discours, de sensibilisation et de dispositifs qui permettent d'apprécier ce que des documents, on va dire, originaux et des documents produits par l'IA peuvent avoir de différence.

\textbf{MF :} C’est un peu lié d'une autre façon aussi au dépôt légal du web qui, du coup, en contient le plus. Il y a une surveillance, justement, de ne pas aspirer de l'IA. 

\textbf{AH :} C'est un vrai enjeu de savoir l'identifier, de savoir quels discours elle porte et par quels moyens elle porte ces discours-là.  Pourquoi est-ce qu'elle cherche à reproduire des codes ? qui sont ceux de la photographie, ou qui sont ceux du reportage, ou qui sont ceux de l'archive en noir et blanc, etc. C'est-à-dire, quelle symbolique, ou quelle portée, ou quel type de discours ils mettent derrière ces dispositifs, qui sont des dispositifs visuels, ou des dispositifs de montage, ou d'animation, etc. Donc, il y a un vrai travail à faire là-dessus, et un vrai travail à faire sur ce qu'est un document audiovisuel. 

\textbf{MF :} Donc c'est une évolution un peu de tout finalement. C'est des chercheurs, des documentalistes et des archives. 

\textbf{AH :} Oui, tout à fait.  Et nous, comme on a un rôle de médiation auprès du public à l’INAthèque et à MPS aussi, c'est quelque chose qu'on doit absolument travailler. C'est vraiment un sujet qui va être important, d'autant qu'on doit aussi beaucoup composer avec l'évolution des méthodes, des manières de consommer la production du visuel. C'est-à-dire qu'aujourd'hui, elle est énormément consommée par extrait, sur des réseaux sociaux, auprès de comptes ou de diffuseurs qui font énormément d'éditorialisation et de curation, comme l’INA par ailleurs.  Et du coup, c'est vrai que des médias comme la télé ou la radio sont peut-être, je ne vais pas faire de généralité, mais ils sont peut-être, en tout cas les jeunes générations d'étudiants, ont peut-être moins l'expérience ou pensent en avoir moins l'expérience que des générations plus anciennes.  Et donc, de la même manière, l'histoire de la technique et l'histoire de ces modes de diffusion sont un petit peu moins connu aujourd'hui, il y a besoin de plus insister dessus.  Et en même temps, nous, on a aussi besoin de mieux connaître, et c'est l'un des enjeux du dépôt légal du web, de mieux connaître les nouveaux modes de consommation de ces contenus audiovisuels aussi, pour adapter nos discours. Parce que c'est via les réseaux sociaux que tout ce qui relève de l'IA se diffuse, des IA génératives se diffusent très largement. Donc là, il y a une forme de rupture dans les modes de consommation et donc dans la connaissance liée à ces médias traditionnels ou à ces techniques traditionnelles de diffusion qu'il faut absolument que nous, on garde en tête.

\section*{Les profils de recherche }

\subsection*{Documentaliste} 

\textbf{MF :} Mais est-ce que du côté des chercheurs, étant donné que récemment, on a eu vraiment un changement rapide avec l'IA, est-ce qu'ils sont demandeurs ou est-ce qu'ils posent des questions en rapport à l'IA 

\textbf{Doc :} C'est une bonne question. Il faut savoir qu'à l’Inathèque, il y a aussi un service qui s'appelle le Lab, qui répond à des demandes de certains chercheurs, souvent plutôt en thèse, des doctorants, ou d’autres projets de recherche, pour des travaux au long cours sur des grosses masses de données. Donc, le Lab répond à ces nouvelles demandes, parce qu'il développe justement des outils spécifiques. Pour nous, au quotidien, la plupart des usagers que l'on reçoit viennent pour des recherches un peu ponctuelles, ou sont des étudiants qui apprennent, qui sont en Master 1, en Master 2, mais qui sont encore dans l'apprentissage du travail de chercheur. Ils travaillent sur des sujets souvent un peu plus délimités, on leur montre les outils. J'étais quand même assez agréablement surprise, parce que je prends quand même des précautions oratoires quand je présente Hyperbase qui a 30 ans d’âge. Et bien, au final, j'ai eu de bons retours d'étudiants qui ont 20 ans, 22. Ils me disaient, oui, c'est logique, c'est clair et tout ça. Et je me suis dit, finalement, ça va. Peut-être que, justement, quand c’est un mode de recherche assez cadré, par rapport à Google ou tout ça, où on se retrouve avec des masses de données qu'on ne sait pas trop gérer... On prend le début, mais après, voilà, la fin, on s'en fiche un peu. Là, au moins, on a un lot résultat qui peut être plus ou moins pertinent, mais qu'on peut compléter, sur lequel on peut rebondir, c'est assez logique, en fait. Et finalement, cette logique-là, c'est peut-être un peu rassurant, même pour les jeunes générations.  Je ne sais pas, il faudrait leur poser la question. Parce que, j'avais un peu peur, évidemment, ces jeunes sont habitués aujourd'hui à faire des recherches avec Copilot, Chat GPT et tout. Hyperbase, ça peut sembler un peu violent. Mais en tout cas, pour ceux à qui j'ai eu affaire, ils n’ont pas eu l’air traumatisés. Il faudrait voir sur la grande majorité. Ceci dit, ceux qui viennent chez nous, ce sont aussi des gens qui s'intéressent, qui sont en histoire, en info-com. Donc, c'est pas tout le monde non plus. Il y a une forme de sélection. Et puis, ils sont amenés par leur professeur, parfois. 

\textbf{MF :} Est-ce qu'ils ont conscience, par exemple, des questions de traitement documentaire à l’INA ? Est-ce qu'ils savent qu'il y a certains documents qui sont traités avec plus de détails que d'autres ? 

\textbf{Doc :} C'est notre rôle justement, ça fait partie des éléments qu'on doit leur donner dès le départ. Ceci dit, s'ils font une recherche un peu ponctuelle, on ne les bassine pas avec trop d'informations. J'essaie de donner uniquement ce dont ils auront besoin. Mais si, par exemple, quelqu'un en Master 2 commence à faire des recherches et travaille sur le vraiment long terme, là, oui, j'avertis tout de suite. Je dis, voilà, attention, vous allez trouver ça, ça va être plus ou moins décrit. Il faudra chercher peut-être d'une autre façon, peut-être travailler plus large ou travailler par des noms de personnalités, pour essayer de récupérer d'autres documents, plutôt que par thématique, ou vice versa, ça peut être selon les thèmes, selon l'époque. Et là, c'est notre connaissance des fonds qui nous permet d'aider. Même si on ne connaît pas tout, parce que nous-mêmes, on en découvre toujours. Mais c'est à nous, justement, de les alerter pour qu'ils ne fassent pas fausse route. Après, notre travail, c'est aussi de les suivre un petit peu, ça c'est plus délicat, ce n'est pas toujours évident, parce qu'ils n'osent pas trop demander. Alors je leur dis qu’il ne faut pas hésiter à venir nous voir, poser des questions. Il y en a c'est vrai qui viennent spontanément. Mais bon, parfois ils se contentent juste de leur petite recherche. Le risque est que parfois on peut effectivement faire fausse route, on se dit : voilà j'ai des résultats, ok c'est bon. C'est comme l'IA, il dit ok je te réponds, c'est peut-être complètement faux, mais c'est bon, je te donne un résultat.  Là aussi, on peut avoir des résultats un peu faussés par rapport à une réalité, la réalité de la diffusion, etc. Par exemple, avant le dépôt légal, on ne conservait à l'INA que ce qui était produit et coproduit par les chaînes. Tout le reste, il n'y a aucune trace par rapport à la diffusion. Donc ça aussi, ça fait partie des choses qu'on doit transmettre. Si quelqu'un veut travailler sur les séries américaines, ça sera après 95, parce qu'avant, on n'en aura pas de traces. Sauf s'il y a un reportage dessus. Enfin voilà, ça fait toutes sortes de petites informations. Et ça, c'est notre travail, justement, à l’INA thèque. C'est d'aider les chercheurs parce que c'est pas forcément intuitif non plus, la recherche, et puis, souvent, les jeunes générations ne connaissent même pas la télévision. Ils ne la regardent plus, en fait, et donc l'histoire de la télévision, de ses personnalités, ça ne leur dit rien, alors que, nous, on a grandi avec. Donc c'est à nous aussi de leur transmettre aussi quelques petites bases de l'histoire de la télé, l'arrivée des chaînes etc. Pourquoi au début il n'y avait qu'une seule chaîne, etc. Ce sont des choses qu'on n'avait pas à raconter avant. Il y a 20 ans, 25 ans, ce n'était pas la peine, on avait une culture un peu commune.  Maintenant, je me rends compte qu'il y a une fracture, quand même.

\textbf{MF :} Même au niveau des chercheurs, il y a une évolution ?   

\textbf{Doc :} On voit bien que le niveau de connaissance, de références, est différent, ça c'est sûr. Par contre, eux, vont me parler des réseaux sociaux. Parce qu'ils peuvent consulter le dépôt légal du web aussi. Je parle toujours de la radio/télé, mais il ne faut pas oublier le dépôt légal du web. C'est quelque chose qui, à mon avis, va prendre de l'importance.  Pour l'instant, on a des outils, on ne décrit pas très finement les choses mais à terme ça va forcément se développer. On voit bien qu'il y a plein de programmes qui ne sont plus du tout diffusés sur le hertzien, en linéaire en tout cas, et des programmes qu'on appelle natifs web. Arte, c'est eux qui ont commencé, c'était parmi les pionniers. Mais aussi France Télé et les autres chaînes ont des programmes qui ne figurent plus que sur le web, qui peuvent être diffusés ensuite, ou non. Donc, il y a le côté linéaire qui, à mon avis, va perdurer encore pour quelques temps. Mais tout le reste, c'est vrai que ça prend une importance. Et puis, ça fait partie du paysage médiatique. C’est aussi la mission du dépôt légal. Ça en fait partie pleinement.

\textbf{MF :} D'où ça ramène au patrimoine. 

\textbf{Doc :} Exactement. Pour avoir une trace. Pour pouvoir retrouver ce programme-là. 

\textbf{MF :} Il y a une évolution du chercheur sur les sujets qu'il recherche ?  

\textbf{Doc :} Oui, je pense, aussi parce que les anciens profs partent à la retraite. Je l'ai observé récemment, depuis deux ans à peu près, trois ans, parce qu'en fait, ce souvent les mêmes professeurs qui envoient les Master 1 chez nous. Les sujets, pendant des années, c'étaient souvent des sujets du genre comparaison du traitement journalistique entre TF1 et France 2 sur une thématique, c'était le truc classique, ou alors parfois des sujets très complexes, mais volontairement. Ce n’étaient pas les étudiants qui choisissaient le sujet, les professeurs leur posaient une question vraiment assez balèze, presque un sujet de thèse, d'une certaine manière, pour voir comment ils allaient se dépatouiller avec l'archive.  Moi, c'est ce que j'ai compris : c'était le but, en fait, de voir comment ils allaient pouvoir s'orienter. Et moi-même, avec ce genre de sujets, je me disais « Qu'est-ce que je fais ? ». Je n'ai pas de sujet précis en tête, il faudrait retrouver les thématiques. Mais là, je trouve que depuis 2-3 ans, ce sont des sujets plus intéressants. De fait, on est plus à l'aise pour répondre, parce que franchement, il y avait des sujets, on ne savait pas bien comment les prendre. Mais c'était fait exprès. C'était une façon de défier les étudiants un petit peu, pour voir comment ils allaient réagir. Bon, c'est une façon de faire.  Moi, je trouvais ça pas forcément si constructif que ça. Alors que maintenant, il y a des sujets que je trouve un peu plus… sur l'histoire des variétés, autour d'une personnalité, enfin quelque chose sur quoi on peut construire une petite stratégie de recherche, montrer comment on élabore le corpus, et justement savoir pourquoi on peut répondre de telle ou telle façon, pourquoi il y aura des choses qu'on n'aura ou pas. Pour un étudiant, je trouve que c'est bien, parce qu'il se confronte aux questions de la recherche. Comme un médiéviste, il va trouver quelques sources, et puis il y a des choses qu'il n'aura pas, puis d'autres qu'il va pouvoir trouver, ou chercher des biais pour contourner la difficulté. C'est exactement pareil avec nos sources. Et ça, je trouve que c'est constructif, ça permet d'apprendre à chercher aussi. Enfin, de mon point de vue. Alors qu'avant, je trouvais que c'était un peu dur quand même, les questions, enfin, les thèmes. 

\textbf{MF :} Là, du coup, on est sur des personnes qui arrivent avec des questions précises qu'on leur a posées, mais est-ce qu'il y a quand même des personnes qui viennent sans idées précises ? 

\textbf{Doc :} Ça arrive, il y a des chercheurs, notamment ceux qui veulent préparer une thèse, qui ont une petite idée, qui cherchent un sujet. Donc parfois, ils arrivent avec deux, trois sujets sans être vraiment décidés. Ils viennent voir un peu, justement, tâter le terrain pour voir ce qu'on peut avoir, comment ça peut répondre, si ça peut être intéressant ou pas, s'il y a de la matière aussi. Donc ça, ça arrive quelques fois, notamment l'été, parce qu'avant la rentrée, il y en a qui viennent en prospective pour savoir s'il y a de la matière, tout simplement, si ça vaut le coup de travailler sur ce lieu là ou pas. 

\subsection*{Pascal Lebeurier }

\textbf{MF :} Que ce soit des différents professionnels, que ce soit l'INA ou des entreprises privées, est-ce que c'est souvent des recherches par thème comme ça ? 

\textbf{PL :} Ça dépend, ça peut être par thème, par personnalité. Quand on travaille, par exemple, pour les missions ADN. Quand on a travaillé sur Nikos Aliagas, par exemple, évidemment, il y avait la Grèce, la dictature des colonels. Il y a aussi des trucs plus conceptuels. Ils nous avaient demandé quelqu'un qui consommait, qui collectionnait tout, de trouver soit quelqu'un qui parle du matérialiste, de vouloir consommer, ou le contraire. Quelqu'un qui dit qu'en fait, il ne faudrait pas posséder de bien matériel. Donc ça peut être plus conceptuel parfois. Et ça peut-être là-dessus que la sémantique va nous aider. Et pour Lumni, parfois, dans des recherches, on a travaillé sur des pastilles philo pour le programme du bac. C'était genre le bonheur, l'Etat. Et donc les recherches parfois étaient assez conceptuelles. Ils voulaient quelqu'un qui explique le bonheur ou quelqu'un qui exprime le bonheur. Et si tu fais la déclaration classique, c'est pas facile. On arrive à trouver mais on met du temps. Mais le conceptuel existe, c'est ça qu'il faut savoir.  Et parfois il faut faire le contraire.  Parfois, pour trouver quelque chose, il faut chercher le contraire. À un moment, on cherchait pour Paris, pour un musée, il y avait une expo sur les clichés parisiens. Ils voulaient qu'on parle de la saleté à Paris, dans les années 60. Et en fait, on l'a trouvé beaucoup par le contraire. En fait, il y avait des campagnes propretés dans les années 60, ça s'appelait « campagne propreté dans Paris ». Et en fait, quand ils parlaient des campagnes, ils parlaient aussi de la saleté. Mais si tu cherchais de la saleté, tu ne le trouvais pas. Donc il fallait chercher le contraire. 

\textbf{MF :} Est-ce que vous ressentez un changement dans les habitudes des professionnels, ou ça reste constant malgré l'évolution ? 

\textbf{PL :} Non, le changement qu'on peut percevoir actuellement, c'est que c'est assez bizarre, c'est que finalement, il y a de plus en plus de demandes d'illustres très précises. Là on travaille sur Céline Dion. Il va y avoir une demande de plans d'illustres énormes. Et parfois ça peut être aussi des illustres de Québec, des illustres face à de l'Olympia, tout ça. Et on sent que de plus en plus, on allait travailler pour un film CNC de 15 minutes, histoire du CNC, quand on regarde la demande, elle faisait 7 pages. Alors c'est ça, t'as l'impression que maintenant, ils construisent très précisément leur demande. Ce qui fait que c'est très chronophage quand tu dois faire la recherche. Là où on travaille pour les productions INA, parfois, les compositions, ça faisait 20 pages. On leur a dit, si, on suit votre guide, on peut appeler ça guide de recherche, on va y passer 30 jours. Donc va falloir peut-être éliminer des trucs mais on sent qu'il y a de plus en plus de demandes hyper précises, c'est ça qui est bizarre mais parce qu'il y a beaucoup maintenant de trucs du montage cut aussi avec du plan par plan et puis comme ils imaginent leurs trucs dans la tête ils font des demandes très précises.  

\textbf{MF :} Dans ce cadre-là, est-ce qu’ils s'imaginent ce que l'INA peut donner, ou est-ce qu'ils connaissent quand même un petit peu les fonds d'où viennent ? 

\textbf{PL :} Non, des fois, ils sont à côté.  Il y a encore des personnes qui nous demandent des trucs années 10, années 20. On va leur dire non, si vous voulez la Première Guerre, vous allez plutôt aller à la CPAD. Si tu veux, les années 30, la crise aux États-Unis dans les années 30, tu vas pas trouver grand-chose à l'INA. Il va falloir plutôt que tu ailles chercher aux Etats-Unis.  Parfois, il y a encore des personnes qui nous demandent des trucs qui ne sont pas à fond INA.  Pour la grande explication, qui est la série pour Lumni, on a travaillé sur « Bloody Sunday », par exemple, tu vas chercher aussi en Angleterre. Et on a fait l'attentat de Sarajevo, c'est pareil, ça ne sera pas à l’INA, il faudra aller chez Gaumont. Donc, non, parfois c'est nous qui leur disons, ben non, ça, ça ne sera pas à l’INA. 

\subsection*{Arthur Hénaff}

\textbf{MF :} Pour revenir sur les habitudes des chercheurs, au niveau de ce qu'ils vont chercher, est-ce qu'ils iront plutôt chercher dans les mots-clés ou dans le tout-texte ?  

\textbf{AH :} Ça, c'est difficile de vous répondre. Et là, c'est plutôt mes collègues responsables documentaires qui pourraient peut-être vous dire, et les documentalistes en salle. Et puis, ça dépend des cas.  Il y a plein de choses qui sont utilisées.  Il y a les descripteurs, évidemment, en mots-clés, qui vont être importants.  Il y a des chercheurs qui vont avoir besoin de retrouver des entités nommées. Il y a des chercheurs qui vont plutôt se concentrer sur les titres.  Il y a des chercheurs pour qui les données intéressantes  ce sera les dates, heures et durées de diffusion. Donc c'est difficile de vous dire, évidemment les recherches, de manière globale, elles se font de toute façon sur l'index général de l'outil de recherche, et puis ensuite c'est affiné, mais il y a plein de champs qui peuvent être sollicités à la fois en recherche, donc pour trouver les documents qu'on veut à travers nos outils de recherche. Mais il y a aussi plein de données différentes qui peuvent être demandées au Lab pour faire des études quanti plus poussées.  Et des données qui sont de plus en plus demandées, et ça c'est vrai que c'est côté Lab qu'on les fournit, mais il y a à mon avis une vraie réflexion à mener sur la possibilité de les mettre en consultation, toutes recherches confondues, c'est les transcriptions. Presque l'intégralité du fonds INA est transcrit et les chercheurs nous disent ça de manière unanime, qu'au-delà même de la question de les récupérer, ce qu'ils voudraient pouvoir faire, c'est de la recherche en plein texte directement dans les transcriptions.  Parce qu'il y a des entités nommées qui vont échapper à l'indexation, parce que l'indexation, c'est des choix, donc il y a des mots qu'on va retenir, d'autres qu'on ne va pas retenir, et qu'eux, ils vont vouloir récupérer. Par exemple, on a ces discussions assez régulièrement. Si on prend des mots comme féminicide, ce sont des mots qui sont apparus assez récemment, qui n'existaient pas jusqu'à un certain temps dans les notices descriptives, dans les notices des documents.  Et donc si on veut faire des recherches sur les féminicides, une recherche en plein texte dans les transcriptions, va faire évidemment ressortir beaucoup plus de résultats qu’une recherche par mot-clé dans les notices. La possibilité de faire des recherches en plein texte, dans les transcriptions speech-to-texte, et un jour de faire aussi des recherches via les transcriptions des images animées, ce seront des choses qui seront très demandées par tous les chercheurs, que ce soit des chercheurs quanti ou quali.

\textbf{MF :} Et est-ce que vous voyez, dans le futur, par rapport aux nouvelles méthodes de recherche, une évolution des profils des chercheurs ? 

\textbf{AH :} Alors, oui. Il y a plusieurs choses, parce que si on parle des outils, il n'y a pas tellement d'évolution de pratiques de recherche documentaire, vu que les outils n'ont pas évolué.  Après, si on parle de pratiques de recherche au sens large, bien sûr, et ça a été toute l'histoire de la création du Lab, qui est une structure assez jeune. Les demandes des chercheurs, elles portent de plus en plus sur des corpus très importants de données et sur de l'analyse quanti, assisté par des outils de traitement automatique du langage, par des outils d'IA, ou en tout cas des approches qui articulent des analyses qualitatives de critiques documentaires et des analyses quantitatives. Et donc ça, la création du Lab, ça a été la réponse de la valorisation scientifique à l’INA, à l'évolution des pratiques des chercheurs. Et c'est pour ça qu'il a une sorte de double rattachement fonctionnel qui est très intéressant, à la fois l'expertise documentaire de l'INA qui s'incarne auprès des chercheurs dans l'INAthèque, et l'expertise aussi en ingénierie qui, elle, s'incarne dans le service recherche de l'INA. C'est pour ça que le lab est copiloté à la fois côté DPAT (direction des patrimoines) et côté DDT.

\textbf{MF :} Questions récurrentes des chercheurs ? 

\textbf{AH :} Les questions récurrentes de la part des chercheurs, donc évidemment, je vous ai dit les transcriptions. Mais ils ont aussi vraiment des questions et un intérêt très vif pour certains, pour la manière dont on produit des métadonnées documentaires, et notamment pour la manière dont on produit des données d'indexation. On a des discussions très intéressantes avec eux, notamment via un dispositif qui est vraiment un dispositif INAthèque/MPS. On travaille main dans la main depuis cette année sur ce sujet-là. Les chercheurs ont des calendriers et des emplois du temps de plus en plus contraints, ils ont moins de temps pour faire de la recherche, ils sont extrêmement sollicités par leurs enseignements, mais aussi par leurs responsabilités administratives et par tout le travail qu'ils doivent faire pour monter des projets de recherche.  Ils ont un petit peu moins de temps, ou en tout cas, c'est des temps plus compliqués, plus contraints pour faire de la recherche.  Et un des besoins qu'ils peuvent avoir, c'est de se retrouver entre eux pour faire des séminaires de recherche où il n'y a pas forcément de public, où il n'y a pas d'attente, où il n'y a pas de publication derrière, tout simplement se retrouver pour travailler au contact des archives, dans un lieu qui met tout à leur disposition, pour qu'ils puissent travailler de manière collective.  Parce que ça, c'est une des évolutions qu'on voit aussi dans les formats pédagogiques, depuis longtemps maintenant. C'est des modes de recherche ou d'apprentissage de plus en plus collectifs et collaboratifs. Et donc, on organise avec MPS et l’INAthèque, en fonction des projets de recherche auxquels on collabore, des ateliers où les chercheurs, sur une thématique (il y avait Boris Vian, il y avait les personnes racisées dans les programmes de variété en France, il y en aura d'autres qui vont venir dès cette année) se retrouvent autour d'une sélection d'archives et discutent entre eux et ils font avancer leurs propres réflexions. Et après, ils partagent leurs réflexions avec nous.  Et on a eu avec certains d'entre eux des discussions vraiment très intéressantes où ils nous disent « voilà comment on a cherché des archives, voilà ce qu'on a noté sur la manière dont c'était décrit ».  Et donc, on a un échange entre d'un côté des chercheurs, de l'autre côté des archivistes, pour le dire de manière très globale, sur ce que pourraient être leurs attentes en termes d'indexation. Et en même temps, sur nous, ce que sont nos objectifs en tant qu'archivistes.  Et ce sont des discussions très intéressantes parce que les chercheurs vont soit élaborer des concepts, soit ils vont redécrire ou renommer, ou nommer pour la première fois des phénomènes d'une certaine manière.  Et ce sont des choses qui peuvent avancer assez vite. Il peut y avoir de nouveaux concepts qui arrivent dans le débat scientifique et dans le débat public, comme « féminicide », ce genre de choses, et qui n'existent pas forcément dans nos notices. Et eux, ils auraient besoin que ce soient des phénomènes qui puissent être décrits dans nos notices documentaires. Et nous, évidemment, on suit ces évolutions, mais nous, ce dont on a besoin, c'est d'avoir des thésaurus qui soient stables et qui permettent de trouver des ensembles assez importants. Et on a besoin aussi de thésaurus et de méthodes de recherche qui soient suffisamment simples pour qu'on puisse les expliquer.  Donc, on ne peut pas aller sur tous les sujets au rythme de la recherche et au rythme des débats scientifiques, qui ne sont pas toujours réglés d'ailleurs, et où il y a des controverses et des débats qui sont tout à fait sains, sur l'utilisation de tel ou tel terme, mais en même temps on est aussi très preneur des retours des chercheurs sur la manière dont on construit des collections, sur la manière dont on décrit et dont on traite un certain nombre de phénomènes. Et ce sont des discussions qui sont à la fois intéressantes pour ce qu'elles nous apprennent de nos collections, en en ce quelles nous permettent aussi de réfléchir à nos pratiques documentaires.  Donc ça, c'est un sujet assez récurrent avec les chercheurs. 

\textbf{MF :} Les données utilisées par les chercheurs ? :  

\textbf{AH :} Sur les données utilisées, c'est difficile de dire parce qu'il y a plein de choses qui peuvent être pertinentes pour les chercheurs.  Ça va dépendre de leur sujet. L'articulation de recherche quanti à de la recherche quali, c'est le besoin de travailler sur des transcriptions, c'est l'envie de travailler sur des données de plus en plus massives. Ça, c'est vraiment quelque chose qu'on note, notamment chez les chercheurs qui candidatent pour les appels à projets du Lab, parce que chaque année, on accompagne vraiment en très grande proximité six projets, en plus de tous les projets qu'on peut accompagner de manière plus ponctuelle. Et donc, on reçoit beaucoup de candidatures. Et on voit qu'il y a des gens qui veulent travailler sur des corpus de plus en plus vastes, enfin sur des ensembles de données de plus en plus vastes. Et de notre côté, il y a un vrai enjeu aussi de sensibiliser à la manière dont on fait de la recherche sur des jeux de données importants et que ça n'exonère pas de faire un corpus et de réfléchir précisément à quelles données on veut avant de pouvoir les récupérer.  Donc il y a un enjeu là-dessus, sur la massification des données, sur l'élargissement des jeux de données qu'on nous demande.  Il y a des évolutions qui sont liées à ce que je vous disais en matière de travail collaboratif et d'enjeux pédagogiques collaboratifs. Et un point très important, et là, c'est plutôt un point de côté INAthèque, parce que je voyais qu'il y avait une question sur ce que l'IA a comme impact sur les pratiques pro. 

\section*{L’outil de cartographie du traitement documentaire } 

\subsection*{Documentaliste} 

\textbf{MF :} Et justement, dans ce cadre...  Est-ce que nous, du coup, le projet de cartographie du traitement documentaire ?  

\textbf{Doc :} Ah oui, à mon avis, oui, parce que ça fait partie, déjà, parce que tu vois, nous, la mémoire, même moi, je ne sais plus. Je me dis telle chaîne, quand ça a dû commencer ? Mais on n'est pas très précis, puis on n'a pas le temps de le chercher.  Et je pense que là, pour le coup, ça donne une trame sérieuse qui permettra de délimiter des périmètres, et de dire, voilà, ça on a, ça on n'a pas, et pour X raisons, c'est ça aussi. Comme tu dis, on en parlait tout à l'heure, c'est que parfois, on ne sait plus pourquoi c'est comme ça. C'est pas toujours facile, quoi.

\textbf{MF :} Oui, puis même pour, des fois, relever des différences qu'on n'a pas relevées, s'il y a des résultats qui ne ressortent pas, on pourrait le percevoir.   

\textbf{Doc :} C'est ça, se référer à ça. Par exemple la chaîne, La 5, parce que pendant un temps, on pouvait pas du tout... maintenant, c'est un peu différent, on peut la consulter, mais il fut un temps, on avait vaguement accès aux notices, mais c'était tout, quoi. Enfin, il y a plein de choses comme ça.  On ne savait plus de quand à quand, parce qu'on a perdu cette mémoire-là.  Et puis le traitement des chaînes du câble satellite, le traitement des chaînes radio, pareil. À partir d’un moment, pourquoi on a arrêté de décrire ça ? Et puis pourquoi, par exemple, à partir du dépôt légal, on décrit le journal du soir et pas les autres éditions du journal ? Enfin, il y a plein de choses comme ça. Le 13h de TF1 par exemple. On a des données, mais en fait, ce sont des choses qui sont importées de TF1... Enfin, tu vois, des détails comme ça. Ça, on n’en parle pas forcément aux étudiants.  Mais c'est vrai que dans un résultat, s'ils veulent faire un corpus, s'ils commencent à vouloir faire des statistiques, là, on dit attention. Parce que c'est là qu'il faut être très, très prudent. Pour le coup, c'est quelque chose qui permettra d'avoir des billes, justement, pour dire pourquoi on a tant de réponses pour telle requête. Parce qu'à cette période-là, il y a des éléments. Les choses qu'on a captées, qu'on n'a pas captées. C'est précieux.

\textbf{MF :} C'est super, c'est l'objectif. Mais dans le cadre des chercheurs, mais en tant que documentaliste, que ce soit...  Pour les autres missions ? 

\textbf{Doc :} Aussi, parce que c'est une trace, c'est notre mémoire, en fait.  Et comme je te dis, même nous, même moi qui ai grandi au DL, je perds le fil. Parfois, je ne sais plus, ou c’est imprécis. C'est pour après aussi. Dans 20 ans, si l'INA existe encore, au moins on aura une trace de ça.  Parce que là, nous, on est encore là. Donc on peut dire, mais oui, ça, oui, pourquoi, ça me revient. Mais au bout d'un moment, on ne sera plus là. Parfois, on peut retrouver des vieux mails. Il n'y a pas longtemps, j'ai retrouvé une info dans un vieux mail de 2006. Mais voilà, ce sont des hasards. J'avais gardé mon mail de 2006 qui traînait au fin fond d'un dossier.  

\subsection*{Pascal Lebeurier }

\textbf{MF :} Dans cette mesure, le projet actuellement sur lequel je travaille, c'est de créer un outil qui permet de créer des visualisations qui permettraient de voir dans le temps quel traitement est appliqué à l'émission. Parce que du coup, ça sous-entend que s'il est traité de cette façon, il y aura un certain type d'indexation. Est-ce que ça, ça peut être intéressant ? 

\textbf{PL :} En fait, ce qui est intéressant, c'est que tu aperçois que maintenant côté FTV, on a des imports assez riches maintenant, avec même l'origine des images, tout ça. Mais plus tu as un import très riche et très plan par plan, plus tu vas avoir du bruit à la recherche. Donc finalement, ça pose d'autres questions. C'est pour ça que ce qui est important, c'est de pouvoir chercher dans des champs différents, c'est d'avoir un champ séquence où tu as uniquement des séquences. Le champ descripteur est capital. Finalement, si tu veux éviter le bruit, tu vas te dire je vais chercher le descripteur. Si le descripteur image est bien mis, si tu cherches des plans de l'Élysée, si tu as un descripteur image, qui est bien mis pour l'Élysée, tu devrais pouvoir t'en sortir avec un descripteur image. Je veux l'Élysée, donc je vais chercher un descripteur image.  Si tu cherches des trucs dans tout texte, tu vas avoir un bruit énorme. C'est pour ça que c'est les deux. Quand tu mets beaucoup, tu vas aussi avoir beaucoup de bruit. Il ne faut pas réfléchir à ce silence. Tu vois, le fait où il n'y a pas assez de choses, il y a aussi quand il y en a trop. Ou alors, dans ce cas-là, il faut qu'il y ait des champs différents. Mais ça, c'est important. 

\textbf{MF :} Donc oui, c'est une double mesure, mais comme ça permettrait de voir justement là où il y a du bruit et là où il y en a moins, ça peut être intéressant ? 

\textbf{PL :} Oui, ça peut être intéressant. Récemment, je pense qu'on a des trucs assez riches. Il y a une époque, c'est pareil, on faisait une annexation un peu trop longue en sport et tout ça.  Finalement, on se retrouve avec des résumés énormes, on mettait tout et ça créait beaucoup de bruit aussi. À l'époque, quand je suis arrivé à l’INA, on faisait des matchs de rugby, on faisait des trucs de sept pages. On mettait des prises en touche. Et finalement, après, heureusement, maintenant, on fait des trucs plus synthétiques. Tour de France, c'est pareil. Maintenant, on va faire un chapeau. L'importance d'avoir un chapeau avec quelle étape c'est, qui a remporté l'étape, qui est maillot jaune, après mais il y a moins de d’écritures de plans en dessous. En fait ce qu'il faudrait : Michael travaille aux frontières, il va dire oui il y a un chapeau pour l'étape après s'il se passe un événement dans l'étape une chute comme ça, faudrait pouvoir le notifier dans les plans parce que c'est ça qu'on va demander après. Mais c'est intéressant parce que tu vois tu es censé dire, je vais faire un chapeau pour tout ce qui est, mais par contre s'il s'est passé un truc, il faudrait que je le mette dans le champ-séquence.  Mais c'est là-dessus comme ça qu'il faudrait réfléchir. Mais ça dépend du nombre. Moi, j'ai fait à la chaîne des étapes du tour de France, de différentes années. Si tu veux aller vite, tu ne peux pas te visionner toutes l’étape.  Donc, tu vas faire un chapeau. Tu vas lire. En même temps, ce qui est intéressant, c'est que tu vas regarder le résumé de l'équipe, parce que s'il y a une chute, peut-être que ça sera dans le résumé de l'équipe. Quand tu travailles sur les matchs de l'équipe de France de foot, c'est pareil. Si tu prends le site de la FFF, c'est marqué le buteur, tout ça, ou alors tu as l’intérieur de l'équipe.  Donc peut-être que la commande papier va t'aider, va te dire ce qui s'est passé.  Et comme ça, tu n'as pas besoin de visionner toutes l’émission. C'est pour ça qu'il y a aussi une réflexion. Je dis aux documentalistes qui font des recherches, souvent, de faire des recherches web. Comme ça, déjà, tu as beaucoup de trucs et ça va t'aider à faire ta recherche dans les fonds de l’INA. Après, les deux côtés, ils n’ont pas tous les mêmes méthodes. Il y en a qui vont faire beaucoup de recherches à côté, d'autres plutôt que dans le fond.  Il y a différentes méthodes. Par exemple, dans le document d'accompagnement, je ne sais pas si tu connais l'outil qui s'appelle le PCM, il y a un outil PCM interne où on a accès aux documents d'accompagnement aussi. Tu fais ta recherche, tu as accès à la fin au document, tu peux le visionner. Tu as aussi accès au conducteur, au document d'accompagnement. Donc ça, c'est intéressant. 

\subsection*{Arthur Hénnaff  }

\textbf{MF :} Est-ce que ce serait quelque chose de pertinent pour les chercheurs et les documentalistes, d’avoir une vision sur le traitement documentaire ? 

\textbf{AH :} De toute façon, que ce soit côté Lab, côté Inathèque, c'est un outil de connaissance des collections et de connaissance de l'état du traitement qui est, à mon avis, essentiel.  Et quand il m'avait été présenté pour la première fois, j'avais trouvé ça vraiment passionnant. Et ça, c'est quelque chose qui est une plus-value énorme dans tous les accompagnements qu'on peut faire, que ce soit en quali ou en quanti. La connaissance des manières dont on a traité les collections et de l'état de traitement et de qui fait quoi et en fonction de quelles règles, sur quoi on intervient directement, qu'est-ce qu'on récupère uniquement comme métadonnées d'import depuis les diffuseurs, où on en est dans le traitement du flux, etc. Tous ces éléments de cartographie du traitement, ce sont des éléments de contexte extrêmement importants qui pour moi sont indissociables de la connaissance des fonds. 

\section*{Contact avec le document}

\subsection*{Documentaliste} 

\textbf{MF :} Oui, donc finalement c’est super intéressant parce que ça permet de figurer comment ça change, pourquoi ça change, l'évolution. D'autant plus que du coup, ça fait état du fait que même le métier de documentaliste, il a énormément évolué depuis ses débuts, en tout cas à l’INA. 

\textbf{Doc :} Déjà, quand tu passes de la notice tapée à la main... puis le visionnage en direct, ils prenaient les notes en direct. Les documentalistes ont eu recours parfois aux magnétoscopes dans les années 80, mais avant, non. Déjà ça, ça en dit long sur les résumés qu'ils pouvaient faire. Heureusement, les reportages à l'époque étaient un peu plus lents qu'aujourd'hui. Mais quand même, il fallait avoir une gymnastique d'esprit. Nous, au Dépôt Légal, on a tout de suite eu les supports physiques, donc ça a aidé. On pouvait faire des pauses sur le document pour noter les noms en incrustation. Avant, il fallait tout noter rapidement et quand il y avait des incrustations... parce qu'il n'y en avait pas non plus toujours. Et pourtant, ils arrivaient quand même à faire du plan par plan...

\textbf{MF :} Par rapport au début du métier de documentaliste à aujourd'hui, est-ce que tu estimes que tu es plus en contact avec le document ou moins étant donné que tout n'est pas traité avec la même importance ? 

\textbf{Doc :} Ça dépend de ce que tu appelles le document. Quand j'étais aux commandes en 1993, avant la numérisation, je faisais des recherches, je sélectionnais les notices, je passais la commande par fax au magasin, je commandais parfois un motard, parce que les supports étaient stockés dans des lieux différents et éloignés parfois. En fait, les images je ne les voyais pas, je ne savais pas en fait ce qui sortait. Pour les voir, comme je travaillais surtout sur TF1, je regardais le journal du soir, pour voir « Tiens ils ont pris ça comme image ». Si j'avais le temps, j'allais à la technique, parce que c'est là que passaient les faisceaux, à l'époque on parlait de faisceaux. Et c'est là que je pouvais voir passer les images sélectionnées. Quand c'étaient des supports films, ça partait de Bry. Il y avait une machine pour envoyer les images à la technique rue de Patay. C’était une espèce de nœud en fait. De là, les images étaient distribuées aux chaînes. Donc ça pouvait être soit des supports films, soit des cassettes vidéo. Différentes variétés de cassettes, d'ailleurs. Donc je ne voyais jamais les images sauf si j'avais le temps de passer en technique. Pour les supports film, le monteur film parfois m'appelait pour dire « bah écoute ton film il faut refaire tous les scotchs donc, pour le 13h, ce sera pas possible ». Parce que, pour le film, tu avais plusieurs matériels pour une émission. Parfois, tu avais des copies, mais parfois juste l'original, avec plein de scotchs de montage. Donc si on pouvait, on commandait en priorité la copie mais parfois il n'y avait que l'original et là on savait que...  parfois tu n'avais pas 36000 choix de documents et donc si le monteur devait refaire tous les scotchs, ça prenait du temps, etc. Parfois, le monteur disait, écoute, ta séquence, elle n'y est pas.  Ça, ça arrivait aussi. Comme, à l'époque, c'était un usage purement professionnel, dans un reportage, quand il y avait une séquence qui intéressait un journaliste, il coupait dedans et la reprenait pour une autre émission, ça arrivait quelquefois. Donc, il fallait repérer l'autre émission pour la commander et récupérer la séquence qu'on cherchait. Donc voilà, ça, c'est par rapport à l'histoire de l’archivage. C'est la période avant la numérisation et c'est vrai que depuis la numérisation, qui a quand même pris une vingtaine d'années, on a accès plus facilement au support.

\section*{Différence dépôt légal/archives professionnelles }

\subsection*{Documentaliste}  

\textbf{MF :} Tous les changements, c'était un peu avant le confinement ? 

\textbf{Doc :} Oui, la grosse réorganisation. En fait, moi-même, j'ai un peu perdu le souvenir des dates exactes.  C’était autour du confinement, plus ou moins. On a quitté nos locaux, le dépôt légal était situé dans des bâtiments à Bry 3, de l'autre côté des décors de ce que l'on appelait la SFP, la Société Française de Production. Il y avait des décors de tournage, on était dans des locaux juste derrière. Donc, on était séparés des autres services de l’INA, physiquement. Déjà, il y a eu un rapatriement et donc un rapprochement des équipes ici à Bry 1, puis la réorganisation des services. Ensuite, il y a le confinement.

\textbf{MF :} Oui donc les deux services, avec d'un côté du coup plus l'utilisation de la base de données du DL, et de l'autre côté, c'était plus Totem. 

\textbf{Doc :} Voilà, c'est ça. C'est qu'en fait, pendant des années, c'était vraiment séparé, on n'avait pas accès à Totem.  On pouvait accéder aux notices des archives professionnelles uniquement par l'outil de consultation qui s'appelle Hyperbase, qui s'appelle toujours Hyperbase. Et, pendant longtemps, on n'avait pas accès aux images et aux sons de ces fonds “professionnels”. On avait uniquement accès à ce qui en avait été copié pour les usagers du dépôt légal. La numérisation a facilité évidemment l’accès. Même côté dépôt légal, avant la mise sur serveur des images ou des sons, il fallait commander les cassettes ou les DVD. C'est au cours des 20 dernières années que l'accès a été facilité.\\